% !TeX root = ./theopendix_check.tex
% Detailed derivations for the theoretical framework in M2Thesis.tex.
% This file is intended to be included with % !TeX root = ./theopendix_check.tex
% Detailed derivations for the theoretical framework in M2Thesis.tex.
% This file is intended to be included with % !TeX root = ./theopendix_check.tex
% Detailed derivations for the theoretical framework in M2Thesis.tex.
% This file is intended to be included with % !TeX root = ./theopendix_check.tex
% Detailed derivations for the theoretical framework in M2Thesis.tex.
% This file is intended to be included with \input{theopendix}.

\section{Theoretical Appendix}
\label{app:theory}

This appendix solves the model in Section 2 step by step. I write the algebra in a way that should be readable for a good undergraduate student who is already comfortable with constrained optimization, CES demand systems, and the idea of a markup as price over marginal cost.

The main goals are:
\begin{enumerate}
    \item derive the demand system at each CES nest from first principles;
    \item derive the inverse-demand system because firms compete in quantities;
    \item show carefully how the perceived elasticity is built from the three substitution margins in the model;
    \item derive the Cournot markup formula and the sectoral markup aggregator;
    \item clarify one notation point: the linear expression in market shares is the \emph{inverse markup} $1/\mu_{ij}$, not the markup $\mu_{ij}$ itself.
\end{enumerate}

\subsection{A useful CES template}

Before solving each layer separately, it is useful to record one generic CES result that we will reuse.

Suppose a composite $X$ is built from a collection of goods $\{x_n\}_{n \in \mathcal{N}}$:
\[
X =
\left(
\sum_{n \in \mathcal{N}}
a_n^{\frac{1}{\sigma}} x_n^{\frac{\sigma-1}{\sigma}}
\right)^{\frac{\sigma}{\sigma-1}},
\qquad \sigma > 1.
\]
Given prices $\{p_n\}$, the dual expenditure-minimization problem is
\[
\min_{\{x_n\}} \sum_{n \in \mathcal{N}} p_n x_n
\qquad \text{s.t.} \qquad
X \geq \bar{X}.
\]
The Lagrangian is
\[
\mathcal{L}
=
\sum_{n \in \mathcal{N}} p_n x_n
+
\lambda
\left[
\bar{X}
-
\left(
\sum_{n \in \mathcal{N}}
a_n^{\frac{1}{\sigma}} x_n^{\frac{\sigma-1}{\sigma}}
\right)^{\frac{\sigma}{\sigma-1}}
\right].
\]

Let
\[
M \equiv \sum_{n \in \mathcal{N}} a_n^{\frac{1}{\sigma}} x_n^{\frac{\sigma-1}{\sigma}},
\qquad
X = M^{\frac{\sigma}{\sigma-1}}.
\]
Then
\[
\frac{\partial X}{\partial x_n}
=
\frac{\sigma}{\sigma-1}
M^{\frac{1}{\sigma-1}}
\cdot
a_n^{\frac{1}{\sigma}}
\cdot
\frac{\sigma-1}{\sigma} x_n^{-\frac{1}{\sigma}}
=
X^{\frac{1}{\sigma}} a_n^{\frac{1}{\sigma}} x_n^{-\frac{1}{\sigma}}.
\]
Therefore the first-order condition for each $n$ is
\begin{equation}
p_n = \lambda X^{\frac{1}{\sigma}} a_n^{\frac{1}{\sigma}} x_n^{-\frac{1}{\sigma}}.
\label{eq:app_generic_ces_foc}
\end{equation}

Take two goods $n$ and $m$. Dividing the two FOCs gives
\[
\frac{p_n}{p_m}
=
\left( \frac{a_n}{a_m} \right)^{\frac{1}{\sigma}}
\left( \frac{x_n}{x_m} \right)^{-\frac{1}{\sigma}}
\Longrightarrow
\frac{x_n}{x_m}
=
\frac{a_n}{a_m}
\left( \frac{p_n}{p_m} \right)^{-\sigma}.
\]
This is the familiar CES relative-demand equation. Solving completely gives
\begin{equation}
x_n = a_n \left( \frac{p_n}{P_X} \right)^{-\sigma} X,
\label{eq:app_generic_hicksian}
\end{equation}
where the CES price index is
\begin{equation}
P_X =
\left(
\sum_{n \in \mathcal{N}}
a_n p_n^{1-\sigma}
\right)^{\frac{1}{1-\sigma}}.
\label{eq:app_generic_price_index}
\end{equation}

Finally, rearranging \eqref{eq:app_generic_hicksian} gives the inverse-demand form
\begin{equation}
p_n
=
P_X
\left( \frac{x_n}{a_n X} \right)^{-\frac{1}{\sigma}}
=
a_n^{\frac{1}{\sigma}} P_X \left( \frac{x_n}{X} \right)^{-\frac{1}{\sigma}}.
\label{eq:app_generic_inverse_demand}
\end{equation}

We now apply the same steps to each nest in the model.

\subsection{Top nest: allocation across sectors}

The representative household values sectoral composites according to
\[
C =
\left(
\int_0^1
\alpha_j^{\frac{1}{\eta}} Q_j^{\frac{\eta-1}{\eta}} dj
\right)^{\frac{\eta}{\eta-1}},
\qquad
\eta > 1,
\qquad
\int_0^1 \alpha_j dj = 1.
\]

\subsubsection{Expenditure minimization problem}

Fix a target final consumption level $\bar{C}$. The household solves
\[
\min_{\{Q_j\}_{j \in [0,1]}}
\int_0^1 P_j Q_j dj
\qquad \text{s.t.} \qquad
\left(
\int_0^1
\alpha_j^{\frac{1}{\eta}} Q_j^{\frac{\eta-1}{\eta}} dj
\right)^{\frac{\eta}{\eta-1}}
\geq \bar{C}.
\]

The Lagrangian is
\[
\mathcal{L}
=
\int_0^1 P_j Q_j dj
+
\lambda
\left[
\bar{C}
-
\left(
\int_0^1
\alpha_j^{\frac{1}{\eta}} Q_j^{\frac{\eta-1}{\eta}} dj
\right)^{\frac{\eta}{\eta-1}}
\right].
\]

Using the generic CES template with a continuum of goods, the FOC for sector $j$ is
\begin{equation}
P_j = \lambda C^{\frac{1}{\eta}} \alpha_j^{\frac{1}{\eta}} Q_j^{-\frac{1}{\eta}}.
\label{eq:app_top_foc}
\end{equation}
Take two sectors $j$ and $\ell$. Dividing their FOCs gives
\[
\frac{Q_j}{Q_\ell}
=
\frac{\alpha_j}{\alpha_\ell}
\left( \frac{P_j}{P_\ell} \right)^{-\eta}.
\]
Hence sectoral demand is
\begin{equation}
Q_j = \alpha_j \left( \frac{P_j}{P} \right)^{-\eta} C,
\label{eq:app_top_demand}
\end{equation}
where
\begin{equation}
P =
\left(
\int_0^1 \alpha_j P_j^{1-\eta} dj
\right)^{\frac{1}{1-\eta}}
\label{eq:app_top_price_index}
\end{equation}
is the final CES price index.

\subsubsection{Sectoral expenditure}

Multiply \eqref{eq:app_top_demand} by $P_j$:
\[
D_j \equiv P_j Q_j = \alpha_j \left( \frac{P_j}{P} \right)^{1-\eta} P C.
\]
Using $Y = PC$, we obtain
\begin{equation}
D_j = \alpha_j \left( \frac{P_j}{P} \right)^{1-\eta} Y.
\label{eq:app_sector_expenditure}
\end{equation}
This is the sector-level expenditure expression reported in the main text.

\subsubsection{Inverse demand at the sector level}

Rearranging \eqref{eq:app_top_demand} gives
\begin{equation}
P_j = P \, \alpha_j^{\frac{1}{\eta}} \left( \frac{Q_j}{C} \right)^{-\frac{1}{\eta}}.
\label{eq:app_top_inverse}
\end{equation}
Taking logs,
\[
\ln P_j = \ln P + \frac{1}{\eta} \ln \alpha_j - \frac{1}{\eta} \ln Q_j + \frac{1}{\eta} \ln C.
\]
So, holding $P$, $\alpha_j$, and $C$ fixed,
\begin{equation}
\frac{d \ln P_j}{d \ln Q_j} = - \frac{1}{\eta}.
\label{eq:app_top_inverse_elasticity}
\end{equation}
This is the inverse-demand elasticity at the top nest.

\subsection{Middle nest: domestic versus foreign supply inside sector}

Within sector $j$, the sectoral composite is
\[
Q_j =
\left[
(\Lambda_j^D)^{\frac{1}{\nu_j}} (Q_j^D)^{\frac{\nu_j-1}{\nu_j}}
+
(1-\Lambda_j^D)^{\frac{1}{\nu_j}} (Q_j^F)^{\frac{\nu_j-1}{\nu_j}}
\right]^{\frac{\nu_j}{\nu_j-1}},
\qquad \nu_j > 1.
\]

\subsubsection{Expenditure minimization problem}

For a given target $\bar{Q}_j$, minimize expenditure on the domestic and foreign composites:
\[
\min_{Q_j^D,Q_j^F}
P_j^D Q_j^D + P_j^F Q_j^F
\qquad \text{s.t.} \qquad
Q_j \geq \bar{Q}_j.
\]
The Lagrangian is
\[
\mathcal{L}
=
P_j^D Q_j^D + P_j^F Q_j^F
+
\lambda_j
\left[
\bar{Q}_j
-
\left[
(\Lambda_j^D)^{\frac{1}{\nu_j}} (Q_j^D)^{\frac{\nu_j-1}{\nu_j}}
+
(1-\Lambda_j^D)^{\frac{1}{\nu_j}} (Q_j^F)^{\frac{\nu_j-1}{\nu_j}}
\right]^{\frac{\nu_j}{\nu_j-1}}
\right].
\]

The FOCs are
\begin{equation}
P_j^D
=
\lambda_j Q_j^{\frac{1}{\nu_j}}
(\Lambda_j^D)^{\frac{1}{\nu_j}}
(Q_j^D)^{-\frac{1}{\nu_j}},
\label{eq:app_middle_foc_D}
\end{equation}
\begin{equation}
P_j^F
=
\lambda_j Q_j^{\frac{1}{\nu_j}}
(1-\Lambda_j^D)^{\frac{1}{\nu_j}}
(Q_j^F)^{-\frac{1}{\nu_j}}.
\label{eq:app_middle_foc_F}
\end{equation}

Dividing \eqref{eq:app_middle_foc_D} by \eqref{eq:app_middle_foc_F} gives
\[
\frac{Q_j^D}{Q_j^F}
=
\frac{\Lambda_j^D}{1-\Lambda_j^D}
\left( \frac{P_j^D}{P_j^F} \right)^{-\nu_j}.
\]
Therefore the Hicksian demands are
\begin{equation}
Q_j^D = \Lambda_j^D \left( \frac{P_j^D}{P_j} \right)^{-\nu_j} Q_j,
\label{eq:app_middle_demand_D}
\end{equation}
\begin{equation}
Q_j^F = (1-\Lambda_j^D) \left( \frac{P_j^F}{P_j} \right)^{-\nu_j} Q_j,
\label{eq:app_middle_demand_F}
\end{equation}
with price index
\begin{equation}
P_j =
\left[
\Lambda_j^D (P_j^D)^{1-\nu_j}
+
(1-\Lambda_j^D) (P_j^F)^{1-\nu_j}
\right]^{\frac{1}{1-\nu_j}}.
\label{eq:app_middle_price_index}
\end{equation}

\subsubsection{A useful notation point}

The CES weight $\Lambda_j^D$ is not, in general, the same thing as the realized domestic expenditure share. The realized share is
\begin{equation}
\lambda_j^D
\equiv
\frac{P_j^D Q_j^D}{P_j Q_j}
=
\Lambda_j^D \left( \frac{P_j^D}{P_j} \right)^{1-\nu_j}.
\label{eq:app_realized_domestic_share}
\end{equation}
In the main text, I often use $\Lambda_j^D$ as reduced-form shorthand for the domestic expenditure share. For the derivations below, it is cleaner to keep the distinction:
\begin{itemize}
    \item $\Lambda_j^D$: CES weight in preferences;
    \item $\lambda_j^D$: realized domestic expenditure share.
\end{itemize}

\subsubsection{Inverse demand for the domestic composite}

Rearranging \eqref{eq:app_middle_demand_D} gives
\begin{equation}
P_j^D
=
P_j
\left( \frac{Q_j^D}{\Lambda_j^D Q_j} \right)^{-\frac{1}{\nu_j}}.
\label{eq:app_middle_inverse_D}
\end{equation}
Taking logs,
\[
\ln P_j^D
=
\ln P_j
- \frac{1}{\nu_j} \ln Q_j^D
+ \frac{1}{\nu_j} \ln Q_j
+ \frac{1}{\nu_j} \ln \Lambda_j^D.
\]
Hence, when $\Lambda_j^D$ is treated as given,
\begin{equation}
d \ln P_j^D
=
d \ln P_j
- \frac{1}{\nu_j} d \ln Q_j^D
+ \frac{1}{\nu_j} d \ln Q_j.
\label{eq:app_middle_inverse_diff}
\end{equation}

\subsection{Bottom nest: domestic varieties inside the domestic composite}

Within the domestic nest of sector $j$, the domestic composite is
\[
Q_j^D =
\left(
\sum_{i=1}^{N_j}
\xi_i^{\frac{1}{\rho_j}} q_{ij}^{\frac{\rho_j-1}{\rho_j}}
\right)^{\frac{\rho_j}{\rho_j-1}},
\qquad \rho_j > 1.
\]

\subsubsection{Expenditure minimization problem}

For a given target $\bar{Q}_j^D$, minimize spending across domestic varieties:
\[
\min_{\{q_{ij}\}_{i=1}^{N_j}}
\sum_{i=1}^{N_j} p_{ij} q_{ij}
\qquad \text{s.t.} \qquad
Q_j^D \geq \bar{Q}_j^D.
\]
The Lagrangian is
\[
\mathcal{L}
=
\sum_{i=1}^{N_j} p_{ij} q_{ij}
+
\kappa_j
\left[
\bar{Q}_j^D
-
\left(
\sum_{i=1}^{N_j}
\xi_i^{\frac{1}{\rho_j}} q_{ij}^{\frac{\rho_j-1}{\rho_j}}
\right)^{\frac{\rho_j}{\rho_j-1}}
\right].
\]

The FOC for variety $i$ is
\begin{equation}
p_{ij}
=
\kappa_j (Q_j^D)^{\frac{1}{\rho_j}}
\xi_i^{\frac{1}{\rho_j}}
q_{ij}^{-\frac{1}{\rho_j}}.
\label{eq:app_bottom_foc}
\end{equation}
Take two varieties $i$ and $h$ and divide the FOCs:
\[
\frac{q_{ij}}{q_{hj}}
=
\frac{\xi_i}{\xi_h}
\left( \frac{p_{ij}}{p_{hj}} \right)^{-\rho_j}.
\]
Therefore the Hicksian demands are
\begin{equation}
q_{ij}
=
\xi_i \left( \frac{p_{ij}}{P_j^D} \right)^{-\rho_j} Q_j^D,
\label{eq:app_bottom_demand}
\end{equation}
where
\begin{equation}
P_j^D
=
\left(
\sum_{i=1}^{N_j}
\xi_i p_{ij}^{1-\rho_j}
\right)^{\frac{1}{1-\rho_j}}
\label{eq:app_bottom_price_index}
\end{equation}
is the CES price index for the domestic composite.

\subsubsection{Inverse demand for a domestic variety}

Rearranging \eqref{eq:app_bottom_demand},
\begin{equation}
p_{ij}
=
P_j^D
\left( \frac{q_{ij}}{\xi_i Q_j^D} \right)^{-\frac{1}{\rho_j}}
=
\xi_i^{\frac{1}{\rho_j}} P_j^D \left( \frac{q_{ij}}{Q_j^D} \right)^{-\frac{1}{\rho_j}}.
\label{eq:app_bottom_inverse}
\end{equation}
Taking logs gives
\begin{equation}
\ln p_{ij}
=
\ln P_j^D
- \frac{1}{\rho_j} \ln q_{ij}
+ \frac{1}{\rho_j} \ln Q_j^D
+ \frac{1}{\rho_j} \ln \xi_i.
\label{eq:app_bottom_inverse_logs}
\end{equation}

\subsubsection{Revenue shares}

Multiply \eqref{eq:app_bottom_demand} by $p_{ij}$ and divide by total domestic expenditure $P_j^D Q_j^D$:
\begin{equation}
\tilde{s}_{ij}
\equiv
\frac{p_{ij} q_{ij}}{P_j^D Q_j^D}
=
\xi_i \left( \frac{p_{ij}}{P_j^D} \right)^{1-\rho_j}.
\label{eq:app_within_share}
\end{equation}
This is firm $i$'s expenditure share within the \emph{domestic} segment of sector $j$.

Its share in \emph{total} sectoral expenditure is
\begin{equation}
s_{ij}
\equiv
\frac{p_{ij} q_{ij}}{P_j Q_j}
=
\frac{P_j^D Q_j^D}{P_j Q_j}
\cdot
\frac{p_{ij} q_{ij}}{P_j^D Q_j^D}
=
\lambda_j^D \tilde{s}_{ij}.
\label{eq:app_total_share}
\end{equation}
If one uses the reduced-form notation of the main text that absorbs relative-price terms into the domestic share, \eqref{eq:app_total_share} becomes $s_{ij} = \Lambda_j^D \tilde{s}_{ij}$.

\subsection{Useful derivative identities}

The next step is to compute how an increase in firm $i$'s own output $q_{ij}$ affects the relevant aggregates $Q_j^D$ and $Q_j$. These derivatives are what allow the three CES nests to show up inside the perceived elasticity.

\subsubsection{How a change in $q_{ij}$ affects the domestic composite $Q_j^D$}

Differentiate the domestic aggregator with respect to $q_{ij}$:
\[
\frac{\partial Q_j^D}{\partial q_{ij}}
=
(Q_j^D)^{\frac{1}{\rho_j}}
\xi_i^{\frac{1}{\rho_j}}
q_{ij}^{-\frac{1}{\rho_j}}.
\]
Multiply by $q_{ij}/Q_j^D$:
\[
\frac{\partial \ln Q_j^D}{\partial \ln q_{ij}}
=
\frac{q_{ij}}{Q_j^D}
\frac{\partial Q_j^D}{\partial q_{ij}}
=
\xi_i^{\frac{1}{\rho_j}} q_{ij}^{\frac{\rho_j-1}{\rho_j}} (Q_j^D)^{-\frac{\rho_j-1}{\rho_j}}.
\]
Now use the inverse-demand equation \eqref{eq:app_bottom_inverse}:
\[
\frac{p_{ij} q_{ij}}{P_j^D Q_j^D}
=
\xi_i^{\frac{1}{\rho_j}} q_{ij}^{\frac{\rho_j-1}{\rho_j}} (Q_j^D)^{-\frac{\rho_j-1}{\rho_j}}
=
\tilde{s}_{ij}.
\]
Therefore
\begin{equation}
\frac{\partial \ln Q_j^D}{\partial \ln q_{ij}} = \tilde{s}_{ij}.
\label{eq:app_dlogQD_dlogq}
\end{equation}
Intuition: if firm $i$ is large within the domestic segment, increasing its own output has a larger effect on the domestic composite.

\subsubsection{How a change in $Q_j^D$ affects the sectoral composite $Q_j$}

Differentiate the middle CES aggregator with respect to $Q_j^D$:
\[
\frac{\partial Q_j}{\partial Q_j^D}
=
Q_j^{\frac{1}{\nu_j}}
(\Lambda_j^D)^{\frac{1}{\nu_j}}
(Q_j^D)^{-\frac{1}{\nu_j}}.
\]
Multiply by $Q_j^D/Q_j$:
\[
\frac{\partial \ln Q_j}{\partial \ln Q_j^D}
=
(\Lambda_j^D)^{\frac{1}{\nu_j}} (Q_j^D)^{\frac{\nu_j-1}{\nu_j}} Q_j^{-\frac{\nu_j-1}{\nu_j}}.
\]
Using \eqref{eq:app_middle_inverse_D}, this expression is exactly the realized domestic expenditure share:
\begin{equation}
\frac{\partial \ln Q_j}{\partial \ln Q_j^D}
=
\frac{P_j^D Q_j^D}{P_j Q_j}
=
\lambda_j^D.
\label{eq:app_dlogQ_dlogQD}
\end{equation}

Combine \eqref{eq:app_dlogQD_dlogq} and \eqref{eq:app_dlogQ_dlogQD}:
\begin{equation}
\frac{\partial \ln Q_j}{\partial \ln q_{ij}}
=
\frac{\partial \ln Q_j}{\partial \ln Q_j^D}
\cdot
\frac{\partial \ln Q_j^D}{\partial \ln q_{ij}}
=
\lambda_j^D \tilde{s}_{ij}
=
s_{ij}.
\label{eq:app_dlogQ_dlogq}
\end{equation}
So the effect of firm $i$'s quantity on the whole sectoral composite is exactly its total sectoral expenditure share.

\subsection{Cost minimization and marginal cost}

Domestic firm $i$ in sector $j$ produces with
\[
q_{ij} = \varphi_{ij} L_{ij}^{\beta_j} M_{ij}^{1-\beta_j},
\qquad 0 < \beta_j < 1.
\]
Given factor prices $w$ and $P_j^M$, the cost-minimization problem is
\[
\min_{L_{ij},M_{ij}}
w L_{ij} + P_j^M M_{ij}
\qquad \text{s.t.} \qquad
q_{ij} = \varphi_{ij} L_{ij}^{\beta_j} M_{ij}^{1-\beta_j}.
\]
The Lagrangian is
\[
\mathcal{L}
=
w L_{ij} + P_j^M M_{ij}
+
\lambda_{ij}
\left[
q_{ij} - \varphi_{ij} L_{ij}^{\beta_j} M_{ij}^{1-\beta_j}
\right].
\]

The FOCs are
\begin{equation}
w
=
\lambda_{ij} \varphi_{ij} \beta_j L_{ij}^{\beta_j-1} M_{ij}^{1-\beta_j},
\label{eq:app_cost_foc_L}
\end{equation}
\begin{equation}
P_j^M
=
\lambda_{ij} \varphi_{ij} (1-\beta_j) L_{ij}^{\beta_j} M_{ij}^{-\beta_j}.
\label{eq:app_cost_foc_M}
\end{equation}

Divide \eqref{eq:app_cost_foc_L} by \eqref{eq:app_cost_foc_M}:
\[
\frac{w}{P_j^M}
=
\frac{\beta_j}{1-\beta_j}
\frac{M_{ij}}{L_{ij}}.
\]
Hence the cost-minimizing input ratio is
\begin{equation}
\frac{M_{ij}}{L_{ij}}
=
\frac{1-\beta_j}{\beta_j} \frac{w}{P_j^M}.
\label{eq:app_input_ratio}
\end{equation}

Now multiply \eqref{eq:app_cost_foc_L} by $L_{ij}$ and \eqref{eq:app_cost_foc_M} by $M_{ij}$:
\[
w L_{ij}
=
\lambda_{ij} \beta_j q_{ij},
\qquad
P_j^M M_{ij}
=
\lambda_{ij} (1-\beta_j) q_{ij}.
\]
Adding the two equations gives
\[
w L_{ij} + P_j^M M_{ij} = \lambda_{ij} q_{ij}.
\]
Because production is constant returns to scale, total cost is proportional to output, so the multiplier $\lambda_{ij}$ is the nominal marginal cost:
\[
\lambda_{ij} = \text{mc}_{ij}.
\]

Solving for the exact unit-cost function yields
\begin{equation}
\text{mc}_{ij}
=
\frac{1}{\varphi_{ij}}
\left( \frac{w}{\beta_j} \right)^{\beta_j}
\left( \frac{P_j^M}{1-\beta_j} \right)^{1-\beta_j}.
\label{eq:app_exact_mc}
\end{equation}
Equivalently,
\[
\text{mc}_{ij}
=
\frac{w^{\beta_j} (P_j^M)^{1-\beta_j}}
{\varphi_{ij} \beta_j^{\beta_j} (1-\beta_j)^{1-\beta_j}}.
\]
Ignoring the sector-specific constant $\beta_j^{-\beta_j} (1-\beta_j)^{-(1-\beta_j)}$, the simpler proportionality used in the main text is
\begin{equation}
\text{mc}_{ij} \propto \varphi_{ij}^{-1} w^{\beta_j} (P_j^M)^{1-\beta_j}.
\label{eq:app_mc_level}
\end{equation}

Taking logs of \eqref{eq:app_exact_mc} or \eqref{eq:app_mc_level} gives
\begin{equation}
d \ln \text{mc}_{ij}
=
- d \ln \varphi_{ij}
+ \beta_j d \ln w
+ (1-\beta_j) d \ln P_j^M.
\label{eq:app_mc_growth}
\end{equation}

\subsection{Input-output propagation of Chinese input shocks}

Let $\mathbf{p}^M$ be the vector of sectoral log input prices, and let $A$ denote the domestic input-coefficient matrix. A first-order approximation of the network structure is
\[
\mathbf{p}^M = A \mathbf{p}^M + \Omega^{CHN} \mathbf{p}^{CHN},
\]
where $\mathbf{p}^{CHN}$ is the vector of Chinese input prices and $\Omega^{CHN}$ collects Chinese sourcing weights.

Rearranging,
\[
(I-A)\mathbf{p}^M = \Omega^{CHN} \mathbf{p}^{CHN},
\qquad
\mathbf{p}^M = (I-A)^{-1} \Omega^{CHN} \mathbf{p}^{CHN}.
\]
Let $L \equiv (I-A)^{-1}$ be the Leontief inverse. For sector $j$,
\begin{equation}
d \ln P_j^M
\approx
\sum_k l_{jk} \omega_k^{CHN} d \ln P_k^{CHN},
\label{eq:app_network_pass}
\end{equation}
where $l_{jk}$ is the $(j,k)$ entry of $L$ and $\omega_k^{CHN}$ is the relevant Chinese sourcing weight.

Substituting \eqref{eq:app_network_pass} into \eqref{eq:app_mc_growth} gives
\begin{equation}
d \ln \text{mc}_{ij}
\approx
- d \ln \varphi_{ij}
+ \beta_j d \ln w
+ (1-\beta_j) \sum_k l_{jk} \omega_k^{CHN} d \ln P_k^{CHN}.
\label{eq:app_mc_reduced}
\end{equation}

\subsection{Deriving the perceived elasticity step by step}

This is the core of the model. Since firms compete in quantities, the cleanest route is:
\begin{enumerate}
    \item write the inverse demand of firm $i$;
    \item differentiate it with respect to $q_{ij}$;
    \item track how a change in $q_{ij}$ moves $Q_j^D$, then $Q_j$, then $P_j$.
\end{enumerate}

\subsubsection{Step 1: inverse demand at the bottom nest}

From \eqref{eq:app_bottom_inverse_logs},
\[
\ln p_{ij}
=
\ln P_j^D
- \frac{1}{\rho_j} \ln q_{ij}
+ \frac{1}{\rho_j} \ln Q_j^D
+ \frac{1}{\rho_j} \ln \xi_i.
\]
Differentiate with respect to $\ln q_{ij}$:
\begin{equation}
\frac{d \ln p_{ij}}{d \ln q_{ij}}
=
\frac{d \ln P_j^D}{d \ln q_{ij}}
- \frac{1}{\rho_j}
+ \frac{1}{\rho_j} \frac{d \ln Q_j^D}{d \ln q_{ij}}.
\label{eq:app_step1}
\end{equation}
Using \eqref{eq:app_dlogQD_dlogq},
\begin{equation}
\frac{d \ln p_{ij}}{d \ln q_{ij}}
=
\frac{d \ln P_j^D}{d \ln q_{ij}}
- \frac{1-\tilde{s}_{ij}}{\rho_j}.
\label{eq:app_step1_simplified}
\end{equation}

\subsubsection{Step 2: how $P_j^D$ responds to $q_{ij}$}

From \eqref{eq:app_middle_inverse_diff},
\[
d \ln P_j^D
=
d \ln P_j
- \frac{1}{\nu_j} d \ln Q_j^D
+ \frac{1}{\nu_j} d \ln Q_j.
\]
Divide by $d \ln q_{ij}$:
\begin{equation}
\frac{d \ln P_j^D}{d \ln q_{ij}}
=
\frac{d \ln P_j}{d \ln q_{ij}}
- \frac{1}{\nu_j} \frac{d \ln Q_j^D}{d \ln q_{ij}}
+ \frac{1}{\nu_j} \frac{d \ln Q_j}{d \ln q_{ij}}.
\label{eq:app_step2}
\end{equation}
Using \eqref{eq:app_dlogQD_dlogq} and \eqref{eq:app_dlogQ_dlogq},
\begin{equation}
\frac{d \ln P_j^D}{d \ln q_{ij}}
=
\frac{d \ln P_j}{d \ln q_{ij}}
- \frac{\tilde{s}_{ij}}{\nu_j}
+ \frac{s_{ij}}{\nu_j}
=
\frac{d \ln P_j}{d \ln q_{ij}}
- \frac{\tilde{s}_{ij} - s_{ij}}{\nu_j}.
\label{eq:app_step2_simplified}
\end{equation}

\subsubsection{Step 3: how $P_j$ responds to $q_{ij}$}

From the sector-level inverse demand \eqref{eq:app_top_inverse},
\[
\ln P_j
=
\ln P
+ \frac{1}{\eta} \ln \alpha_j
- \frac{1}{\eta} \ln Q_j
+ \frac{1}{\eta} \ln C.
\]
Holding the aggregate objects $P$ and $C$ fixed from the viewpoint of an individual firm inside sector $j$,
\begin{equation}
\frac{d \ln P_j}{d \ln q_{ij}}
=
- \frac{1}{\eta} \frac{d \ln Q_j}{d \ln q_{ij}}
=
- \frac{s_{ij}}{\eta},
\label{eq:app_step3}
\end{equation}
where the last equality uses \eqref{eq:app_dlogQ_dlogq}.

\subsubsection{Step 4: combine the three pieces}

Insert \eqref{eq:app_step3} into \eqref{eq:app_step2_simplified}:
\[
\frac{d \ln P_j^D}{d \ln q_{ij}}
=
- \frac{\tilde{s}_{ij} - s_{ij}}{\nu_j}
- \frac{s_{ij}}{\eta}.
\]
Now plug this into \eqref{eq:app_step1_simplified}:
\[
\frac{d \ln p_{ij}}{d \ln q_{ij}}
=
- \frac{1-\tilde{s}_{ij}}{\rho_j}
- \frac{\tilde{s}_{ij} - s_{ij}}{\nu_j}
- \frac{s_{ij}}{\eta}.
\]
Therefore
\begin{equation}
- \frac{d \ln p_{ij}}{d \ln q_{ij}}
=
\frac{1-\tilde{s}_{ij}}{\rho_j}
+
\frac{\tilde{s}_{ij} - s_{ij}}{\nu_j}
+
\frac{s_{ij}}{\eta}.
\label{eq:app_inverse_demand_elasticity}
\end{equation}

Define the firm's \emph{perceived elasticity} under Cournot competition as
\begin{equation}
\frac{1}{\sigma_{ij}^C}
\equiv
- \frac{d \ln p_{ij}}{d \ln q_{ij}}.
\label{eq:app_sigma_definition}
\end{equation}
Equivalently, since price and quantity are locally inverse functions,
\[
\sigma_{ij}^C
=
- \frac{d \ln q_{ij}}{d \ln p_{ij}}.
\]
Combining \eqref{eq:app_inverse_demand_elasticity} and \eqref{eq:app_sigma_definition} yields
\begin{equation}
\sigma_{ij}^C
=
\left(
\frac{1-\tilde{s}_{ij}}{\rho_j}
+
\frac{\tilde{s}_{ij} - s_{ij}}{\nu_j}
+
\frac{s_{ij}}{\eta}
\right)^{-1}.
\label{eq:app_sigma_nested}
\end{equation}

\paragraph{Interpretation.}
Equation \eqref{eq:app_sigma_nested} is easiest to read term by term.
\begin{itemize}
    \item The term $\frac{1-\tilde{s}_{ij}}{\rho_j}$ comes from substitution across \emph{domestic varieties}. If the firm is small within the domestic segment, this force is strong.
    \item The term $\frac{\tilde{s}_{ij}-s_{ij}}{\nu_j}$ comes from substitution between the \emph{domestic composite} and the \emph{foreign composite} inside the sector.
    \item The term $\frac{s_{ij}}{\eta}$ comes from substitution across \emph{sectors} in final demand.
\end{itemize}
Because $\rho_j > \nu_j > \eta$, the inverse elasticities satisfy
\[
\frac{1}{\rho_j} < \frac{1}{\nu_j} < \frac{1}{\eta}.
\]
So as a firm becomes more important in broader and broader aggregates, the relevant residual demand it faces becomes less elastic.

\subsection{Cournot firm's FOC and the markup formula}

Firm $i$ chooses quantity $q_{ij}$ to maximize
\[
\pi_{ij}
=
\bigl(p_{ij}(q_{ij},q_{-ij}) - \text{mc}_{ij}\bigr) q_{ij} - f_{ij}.
\]
The first-order condition is
\begin{equation}
\frac{\partial \pi_{ij}}{\partial q_{ij}}
=
p_{ij}
+ q_{ij} \frac{\partial p_{ij}}{\partial q_{ij}}
- \text{mc}_{ij}
=
0.
\label{eq:app_cournot_foc}
\end{equation}

Divide \eqref{eq:app_cournot_foc} by $p_{ij}$:
\[
1
+ \frac{q_{ij}}{p_{ij}} \frac{\partial p_{ij}}{\partial q_{ij}}
- \frac{\text{mc}_{ij}}{p_{ij}}
=
0.
\]
Using the definition of the perceived elasticity,
\[
\frac{q_{ij}}{p_{ij}} \frac{\partial p_{ij}}{\partial q_{ij}}
=
\frac{d \ln p_{ij}}{d \ln q_{ij}}
=
- \frac{1}{\sigma_{ij}^C}.
\]
Hence
\[
1 - \frac{1}{\sigma_{ij}^C} = \frac{\text{mc}_{ij}}{p_{ij}}.
\]
Define the markup
\[
\mu_{ij} \equiv \frac{p_{ij}}{\text{mc}_{ij}}.
\]
Then
\begin{equation}
\frac{1}{\mu_{ij}}
=
1 - \frac{1}{\sigma_{ij}^C},
\qquad
\mu_{ij}
=
\frac{1}{1-\frac{1}{\sigma_{ij}^C}}
=
\frac{\sigma_{ij}^C}{\sigma_{ij}^C-1}.
\label{eq:app_markup_sigma}
\end{equation}

Substituting \eqref{eq:app_sigma_nested} into \eqref{eq:app_markup_sigma},
\begin{equation}
\frac{1}{\mu_{ij}}
=
1
-
\left[
\frac{1-\tilde{s}_{ij}}{\rho_j}
+
\frac{\tilde{s}_{ij} - s_{ij}}{\nu_j}
+
\frac{s_{ij}}{\eta}
\right].
\label{eq:app_inverse_markup_unsimplified}
\end{equation}
This is already a useful formula: a larger firm faces a less elastic residual demand, so $1/\mu_{ij}$ is lower and $\mu_{ij}$ is higher.

Now use $s_{ij} = \lambda_j^D \tilde{s}_{ij}$, so $\tilde{s}_{ij} = s_{ij}/\lambda_j^D$. Then \eqref{eq:app_inverse_markup_unsimplified} becomes
\begin{align}
\frac{1}{\mu_{ij}}
&=
1 - \frac{1}{\rho_j}
+ \frac{\tilde{s}_{ij}}{\rho_j}
- \frac{\tilde{s}_{ij}}{\nu_j}
+ \frac{s_{ij}}{\nu_j}
- \frac{s_{ij}}{\eta}
\notag \\
&=
\left( 1 - \frac{1}{\rho_j} \right)
-
\left[
\frac{1}{\lambda_j^D}
\left( \frac{1}{\nu_j} - \frac{1}{\rho_j} \right)
+
\left( \frac{1}{\eta} - \frac{1}{\nu_j} \right)
\right] s_{ij}.
\label{eq:app_inverse_markup_linear}
\end{align}

\paragraph{Important notation point.}
Equation \eqref{eq:app_inverse_markup_linear} is a linear equation for the \emph{inverse markup} $1/\mu_{ij}$, not for $\mu_{ij}$ itself. Therefore the corresponding markup is
\begin{equation}
\mu_{ij}
=
\frac{1}{
\left( 1 - \frac{1}{\rho_j} \right)
-
\left[
\frac{1}{\lambda_j^D}
\left( \frac{1}{\nu_j} - \frac{1}{\rho_j} \right)
+
\left( \frac{1}{\eta} - \frac{1}{\nu_j} \right)
\right] s_{ij}
}.
\label{eq:app_markup_full}
\end{equation}
In the main text, I suppress the distinction between $\lambda_j^D$ and $\Lambda_j^D$ for expositional simplicity; under that notation, \eqref{eq:app_inverse_markup_linear} becomes exactly the linear markup condition reported there, except that it should strictly be read as an inverse-markup condition.

\subsection{Sectoral markup and concentration}

Define the sectoral markup as the share-weighted harmonic mean of firm markups:
\begin{equation}
\mu_j
\equiv
\left(
\sum_{i=1}^{N_j} \frac{s_{ij}}{\mu_{ij}}
\right)^{-1}.
\label{eq:app_sector_harmonic}
\end{equation}
Taking inverses,
\[
\frac{1}{\mu_j}
=
\sum_{i=1}^{N_j} s_{ij} \frac{1}{\mu_{ij}}.
\]
Now substitute \eqref{eq:app_inverse_markup_linear}:
\begin{align}
\frac{1}{\mu_j}
&=
\sum_{i=1}^{N_j}
s_{ij}
\left[
\left( 1 - \frac{1}{\rho_j} \right)
-
\left(
\frac{1}{\lambda_j^D}
\left( \frac{1}{\nu_j} - \frac{1}{\rho_j} \right)
+
\left( \frac{1}{\eta} - \frac{1}{\nu_j} \right)
\right) s_{ij}
\right]
\notag \\
&=
\left( 1 - \frac{1}{\rho_j} \right) \sum_{i=1}^{N_j} s_{ij}
-
\left[
\frac{1}{\lambda_j^D}
\left( \frac{1}{\nu_j} - \frac{1}{\rho_j} \right)
+
\left( \frac{1}{\eta} - \frac{1}{\nu_j} \right)
\right]
\sum_{i=1}^{N_j} s_{ij}^2.
\label{eq:app_sector_inverse_step}
\end{align}
Since $\sum_i s_{ij} = 1$, and defining the Herfindahl index
\[
\text{HHI}_j \equiv \sum_{i=1}^{N_j} s_{ij}^2,
\]
we obtain
\begin{equation}
\frac{1}{\mu_j}
=
\left( 1 - \frac{1}{\rho_j} \right)
-
\left[
\frac{1}{\lambda_j^D}
\left( \frac{1}{\nu_j} - \frac{1}{\rho_j} \right)
+
\left( \frac{1}{\eta} - \frac{1}{\nu_j} \right)
\right]
\text{HHI}_j.
\label{eq:app_sector_inverse_markup}
\end{equation}
Therefore
\begin{equation}
\mu_j
=
\frac{1}{
\left( 1 - \frac{1}{\rho_j} \right)
-
\left[
\frac{1}{\lambda_j^D}
\left( \frac{1}{\nu_j} - \frac{1}{\rho_j} \right)
+
\left( \frac{1}{\eta} - \frac{1}{\nu_j} \right)
\right]
\text{HHI}_j
}.
\label{eq:app_sector_markup}
\end{equation}

Again, the \emph{linear} object is $1/\mu_j$, not $\mu_j$.

\subsubsection{Why concentration matters}

Let
\[
A_j \equiv 1 - \frac{1}{\rho_j},
\qquad
B_j \equiv
\frac{1}{\lambda_j^D}
\left( \frac{1}{\nu_j} - \frac{1}{\rho_j} \right)
+
\left( \frac{1}{\eta} - \frac{1}{\nu_j} \right).
\]
Then
\[
\frac{1}{\mu_j} = A_j - B_j \text{HHI}_j,
\qquad
\mu_j = \frac{1}{A_j - B_j \text{HHI}_j}.
\]
Differentiating,
\[
\frac{\partial \mu_j}{\partial \text{HHI}_j}
=
\frac{B_j}{(A_j - B_j \text{HHI}_j)^2}.
\]
Under $\rho_j > \nu_j > \eta$, we have
\[
\frac{1}{\nu_j} - \frac{1}{\rho_j} > 0,
\qquad
\frac{1}{\eta} - \frac{1}{\nu_j} > 0,
\]
so $B_j > 0$. Hence
\begin{equation}
\frac{\partial \mu_j}{\partial \text{HHI}_j} > 0.
\label{eq:app_markup_concentration_positive}
\end{equation}
Sectoral markups rise with concentration because more expenditure gets concentrated on firms that face less elastic residual demand.

\subsection{Dynamic extension and endogenous exit}

Let $x_{ijt}$ denote the state of firm $i$ in sector $j$ at date $t$. This state includes productivity, demand shifters, lagged market share, and the output- and input-trade states $(\tau_{jt}^O,\tau_{jt}^I)$.

If the firm stays active, its value function is
\begin{equation}
V(x_{ijt})
=
\max
\left\{
0,
\pi(x_{ijt})
+ \beta \mathbb{E}\left[ V(x_{ij,t+1}) \mid x_{ijt} \right]
\right\},
\label{eq:app_bellman}
\end{equation}
where $\beta \in (0,1)$ is the discount factor.

Define the continuation surplus
\[
W(x_{ijt})
\equiv
\pi(x_{ijt})
+ \beta \mathbb{E}\left[ V(x_{ij,t+1}) \mid x_{ijt} \right].
\]
The optimal exit rule is
\begin{equation}
\text{Exit}_{ijt} = 1
\quad \Longleftrightarrow \quad
W(x_{ijt}) < 0.
\label{eq:app_exit_rule}
\end{equation}
Equivalently, if current profitability can be summarized by a latent scalar $\theta_{ijt}$, there exists a cutoff function $\theta^*(x_{ijt})$ such that
\begin{equation}
\text{Exit}_{ijt} = 1
\quad \Longleftrightarrow \quad
\theta_{ijt} < \theta^*(x_{ijt}).
\label{eq:app_exit_threshold}
\end{equation}

\paragraph{Why this matters empirically.}
Trade shocks affect both current profits and future continuation values. A rise in final-good import competition may reduce current operating profits; a favorable input-supply shock may reduce marginal cost and increase the value of staying active. Therefore the sample of surviving firms is not random. Observed firm-level markup changes mix:
\begin{itemize}
    \item genuine within-incumbent adjustment;
    \item reallocation across surviving incumbents;
    \item selection induced by endogenous exit.
\end{itemize}
This is the reason the empirical section later introduces a selection-correction strategy.
.

\section{Theoretical Appendix}
\label{app:theory}

This appendix solves the model in Section 2 step by step. I write the algebra in a way that should be readable for a good undergraduate student who is already comfortable with constrained optimization, CES demand systems, and the idea of a markup as price over marginal cost.

The main goals are:
\begin{enumerate}
    \item derive the demand system at each CES nest from first principles;
    \item derive the inverse-demand system because firms compete in quantities;
    \item show carefully how the perceived elasticity is built from the three substitution margins in the model;
    \item derive the Cournot markup formula and the sectoral markup aggregator;
    \item clarify one notation point: the linear expression in market shares is the \emph{inverse markup} $1/\mu_{ij}$, not the markup $\mu_{ij}$ itself.
\end{enumerate}

\subsection{A useful CES template}

Before solving each layer separately, it is useful to record one generic CES result that we will reuse.

Suppose a composite $X$ is built from a collection of goods $\{x_n\}_{n \in \mathcal{N}}$:
\[
X =
\left(
\sum_{n \in \mathcal{N}}
a_n^{\frac{1}{\sigma}} x_n^{\frac{\sigma-1}{\sigma}}
\right)^{\frac{\sigma}{\sigma-1}},
\qquad \sigma > 1.
\]
Given prices $\{p_n\}$, the dual expenditure-minimization problem is
\[
\min_{\{x_n\}} \sum_{n \in \mathcal{N}} p_n x_n
\qquad \text{s.t.} \qquad
X \geq \bar{X}.
\]
The Lagrangian is
\[
\mathcal{L}
=
\sum_{n \in \mathcal{N}} p_n x_n
+
\lambda
\left[
\bar{X}
-
\left(
\sum_{n \in \mathcal{N}}
a_n^{\frac{1}{\sigma}} x_n^{\frac{\sigma-1}{\sigma}}
\right)^{\frac{\sigma}{\sigma-1}}
\right].
\]

Let
\[
M \equiv \sum_{n \in \mathcal{N}} a_n^{\frac{1}{\sigma}} x_n^{\frac{\sigma-1}{\sigma}},
\qquad
X = M^{\frac{\sigma}{\sigma-1}}.
\]
Then
\[
\frac{\partial X}{\partial x_n}
=
\frac{\sigma}{\sigma-1}
M^{\frac{1}{\sigma-1}}
\cdot
a_n^{\frac{1}{\sigma}}
\cdot
\frac{\sigma-1}{\sigma} x_n^{-\frac{1}{\sigma}}
=
X^{\frac{1}{\sigma}} a_n^{\frac{1}{\sigma}} x_n^{-\frac{1}{\sigma}}.
\]
Therefore the first-order condition for each $n$ is
\begin{equation}
p_n = \lambda X^{\frac{1}{\sigma}} a_n^{\frac{1}{\sigma}} x_n^{-\frac{1}{\sigma}}.
\label{eq:app_generic_ces_foc}
\end{equation}

Take two goods $n$ and $m$. Dividing the two FOCs gives
\[
\frac{p_n}{p_m}
=
\left( \frac{a_n}{a_m} \right)^{\frac{1}{\sigma}}
\left( \frac{x_n}{x_m} \right)^{-\frac{1}{\sigma}}
\Longrightarrow
\frac{x_n}{x_m}
=
\frac{a_n}{a_m}
\left( \frac{p_n}{p_m} \right)^{-\sigma}.
\]
This is the familiar CES relative-demand equation. Solving completely gives
\begin{equation}
x_n = a_n \left( \frac{p_n}{P_X} \right)^{-\sigma} X,
\label{eq:app_generic_hicksian}
\end{equation}
where the CES price index is
\begin{equation}
P_X =
\left(
\sum_{n \in \mathcal{N}}
a_n p_n^{1-\sigma}
\right)^{\frac{1}{1-\sigma}}.
\label{eq:app_generic_price_index}
\end{equation}

Finally, rearranging \eqref{eq:app_generic_hicksian} gives the inverse-demand form
\begin{equation}
p_n
=
P_X
\left( \frac{x_n}{a_n X} \right)^{-\frac{1}{\sigma}}
=
a_n^{\frac{1}{\sigma}} P_X \left( \frac{x_n}{X} \right)^{-\frac{1}{\sigma}}.
\label{eq:app_generic_inverse_demand}
\end{equation}

We now apply the same steps to each nest in the model.

\subsection{Top nest: allocation across sectors}

The representative household values sectoral composites according to
\[
C =
\left(
\int_0^1
\alpha_j^{\frac{1}{\eta}} Q_j^{\frac{\eta-1}{\eta}} dj
\right)^{\frac{\eta}{\eta-1}},
\qquad
\eta > 1,
\qquad
\int_0^1 \alpha_j dj = 1.
\]

\subsubsection{Expenditure minimization problem}

Fix a target final consumption level $\bar{C}$. The household solves
\[
\min_{\{Q_j\}_{j \in [0,1]}}
\int_0^1 P_j Q_j dj
\qquad \text{s.t.} \qquad
\left(
\int_0^1
\alpha_j^{\frac{1}{\eta}} Q_j^{\frac{\eta-1}{\eta}} dj
\right)^{\frac{\eta}{\eta-1}}
\geq \bar{C}.
\]

The Lagrangian is
\[
\mathcal{L}
=
\int_0^1 P_j Q_j dj
+
\lambda
\left[
\bar{C}
-
\left(
\int_0^1
\alpha_j^{\frac{1}{\eta}} Q_j^{\frac{\eta-1}{\eta}} dj
\right)^{\frac{\eta}{\eta-1}}
\right].
\]

Using the generic CES template with a continuum of goods, the FOC for sector $j$ is
\begin{equation}
P_j = \lambda C^{\frac{1}{\eta}} \alpha_j^{\frac{1}{\eta}} Q_j^{-\frac{1}{\eta}}.
\label{eq:app_top_foc}
\end{equation}
Take two sectors $j$ and $\ell$. Dividing their FOCs gives
\[
\frac{Q_j}{Q_\ell}
=
\frac{\alpha_j}{\alpha_\ell}
\left( \frac{P_j}{P_\ell} \right)^{-\eta}.
\]
Hence sectoral demand is
\begin{equation}
Q_j = \alpha_j \left( \frac{P_j}{P} \right)^{-\eta} C,
\label{eq:app_top_demand}
\end{equation}
where
\begin{equation}
P =
\left(
\int_0^1 \alpha_j P_j^{1-\eta} dj
\right)^{\frac{1}{1-\eta}}
\label{eq:app_top_price_index}
\end{equation}
is the final CES price index.

\subsubsection{Sectoral expenditure}

Multiply \eqref{eq:app_top_demand} by $P_j$:
\[
D_j \equiv P_j Q_j = \alpha_j \left( \frac{P_j}{P} \right)^{1-\eta} P C.
\]
Using $Y = PC$, we obtain
\begin{equation}
D_j = \alpha_j \left( \frac{P_j}{P} \right)^{1-\eta} Y.
\label{eq:app_sector_expenditure}
\end{equation}
This is the sector-level expenditure expression reported in the main text.

\subsubsection{Inverse demand at the sector level}

Rearranging \eqref{eq:app_top_demand} gives
\begin{equation}
P_j = P \, \alpha_j^{\frac{1}{\eta}} \left( \frac{Q_j}{C} \right)^{-\frac{1}{\eta}}.
\label{eq:app_top_inverse}
\end{equation}
Taking logs,
\[
\ln P_j = \ln P + \frac{1}{\eta} \ln \alpha_j - \frac{1}{\eta} \ln Q_j + \frac{1}{\eta} \ln C.
\]
So, holding $P$, $\alpha_j$, and $C$ fixed,
\begin{equation}
\frac{d \ln P_j}{d \ln Q_j} = - \frac{1}{\eta}.
\label{eq:app_top_inverse_elasticity}
\end{equation}
This is the inverse-demand elasticity at the top nest.

\subsection{Middle nest: domestic versus foreign supply inside sector}

Within sector $j$, the sectoral composite is
\[
Q_j =
\left[
(\Lambda_j^D)^{\frac{1}{\nu_j}} (Q_j^D)^{\frac{\nu_j-1}{\nu_j}}
+
(1-\Lambda_j^D)^{\frac{1}{\nu_j}} (Q_j^F)^{\frac{\nu_j-1}{\nu_j}}
\right]^{\frac{\nu_j}{\nu_j-1}},
\qquad \nu_j > 1.
\]

\subsubsection{Expenditure minimization problem}

For a given target $\bar{Q}_j$, minimize expenditure on the domestic and foreign composites:
\[
\min_{Q_j^D,Q_j^F}
P_j^D Q_j^D + P_j^F Q_j^F
\qquad \text{s.t.} \qquad
Q_j \geq \bar{Q}_j.
\]
The Lagrangian is
\[
\mathcal{L}
=
P_j^D Q_j^D + P_j^F Q_j^F
+
\lambda_j
\left[
\bar{Q}_j
-
\left[
(\Lambda_j^D)^{\frac{1}{\nu_j}} (Q_j^D)^{\frac{\nu_j-1}{\nu_j}}
+
(1-\Lambda_j^D)^{\frac{1}{\nu_j}} (Q_j^F)^{\frac{\nu_j-1}{\nu_j}}
\right]^{\frac{\nu_j}{\nu_j-1}}
\right].
\]

The FOCs are
\begin{equation}
P_j^D
=
\lambda_j Q_j^{\frac{1}{\nu_j}}
(\Lambda_j^D)^{\frac{1}{\nu_j}}
(Q_j^D)^{-\frac{1}{\nu_j}},
\label{eq:app_middle_foc_D}
\end{equation}
\begin{equation}
P_j^F
=
\lambda_j Q_j^{\frac{1}{\nu_j}}
(1-\Lambda_j^D)^{\frac{1}{\nu_j}}
(Q_j^F)^{-\frac{1}{\nu_j}}.
\label{eq:app_middle_foc_F}
\end{equation}

Dividing \eqref{eq:app_middle_foc_D} by \eqref{eq:app_middle_foc_F} gives
\[
\frac{Q_j^D}{Q_j^F}
=
\frac{\Lambda_j^D}{1-\Lambda_j^D}
\left( \frac{P_j^D}{P_j^F} \right)^{-\nu_j}.
\]
Therefore the Hicksian demands are
\begin{equation}
Q_j^D = \Lambda_j^D \left( \frac{P_j^D}{P_j} \right)^{-\nu_j} Q_j,
\label{eq:app_middle_demand_D}
\end{equation}
\begin{equation}
Q_j^F = (1-\Lambda_j^D) \left( \frac{P_j^F}{P_j} \right)^{-\nu_j} Q_j,
\label{eq:app_middle_demand_F}
\end{equation}
with price index
\begin{equation}
P_j =
\left[
\Lambda_j^D (P_j^D)^{1-\nu_j}
+
(1-\Lambda_j^D) (P_j^F)^{1-\nu_j}
\right]^{\frac{1}{1-\nu_j}}.
\label{eq:app_middle_price_index}
\end{equation}

\subsubsection{A useful notation point}

The CES weight $\Lambda_j^D$ is not, in general, the same thing as the realized domestic expenditure share. The realized share is
\begin{equation}
\lambda_j^D
\equiv
\frac{P_j^D Q_j^D}{P_j Q_j}
=
\Lambda_j^D \left( \frac{P_j^D}{P_j} \right)^{1-\nu_j}.
\label{eq:app_realized_domestic_share}
\end{equation}
In the main text, I often use $\Lambda_j^D$ as reduced-form shorthand for the domestic expenditure share. For the derivations below, it is cleaner to keep the distinction:
\begin{itemize}
    \item $\Lambda_j^D$: CES weight in preferences;
    \item $\lambda_j^D$: realized domestic expenditure share.
\end{itemize}

\subsubsection{Inverse demand for the domestic composite}

Rearranging \eqref{eq:app_middle_demand_D} gives
\begin{equation}
P_j^D
=
P_j
\left( \frac{Q_j^D}{\Lambda_j^D Q_j} \right)^{-\frac{1}{\nu_j}}.
\label{eq:app_middle_inverse_D}
\end{equation}
Taking logs,
\[
\ln P_j^D
=
\ln P_j
- \frac{1}{\nu_j} \ln Q_j^D
+ \frac{1}{\nu_j} \ln Q_j
+ \frac{1}{\nu_j} \ln \Lambda_j^D.
\]
Hence, when $\Lambda_j^D$ is treated as given,
\begin{equation}
d \ln P_j^D
=
d \ln P_j
- \frac{1}{\nu_j} d \ln Q_j^D
+ \frac{1}{\nu_j} d \ln Q_j.
\label{eq:app_middle_inverse_diff}
\end{equation}

\subsection{Bottom nest: domestic varieties inside the domestic composite}

Within the domestic nest of sector $j$, the domestic composite is
\[
Q_j^D =
\left(
\sum_{i=1}^{N_j}
\xi_i^{\frac{1}{\rho_j}} q_{ij}^{\frac{\rho_j-1}{\rho_j}}
\right)^{\frac{\rho_j}{\rho_j-1}},
\qquad \rho_j > 1.
\]

\subsubsection{Expenditure minimization problem}

For a given target $\bar{Q}_j^D$, minimize spending across domestic varieties:
\[
\min_{\{q_{ij}\}_{i=1}^{N_j}}
\sum_{i=1}^{N_j} p_{ij} q_{ij}
\qquad \text{s.t.} \qquad
Q_j^D \geq \bar{Q}_j^D.
\]
The Lagrangian is
\[
\mathcal{L}
=
\sum_{i=1}^{N_j} p_{ij} q_{ij}
+
\kappa_j
\left[
\bar{Q}_j^D
-
\left(
\sum_{i=1}^{N_j}
\xi_i^{\frac{1}{\rho_j}} q_{ij}^{\frac{\rho_j-1}{\rho_j}}
\right)^{\frac{\rho_j}{\rho_j-1}}
\right].
\]

The FOC for variety $i$ is
\begin{equation}
p_{ij}
=
\kappa_j (Q_j^D)^{\frac{1}{\rho_j}}
\xi_i^{\frac{1}{\rho_j}}
q_{ij}^{-\frac{1}{\rho_j}}.
\label{eq:app_bottom_foc}
\end{equation}
Take two varieties $i$ and $h$ and divide the FOCs:
\[
\frac{q_{ij}}{q_{hj}}
=
\frac{\xi_i}{\xi_h}
\left( \frac{p_{ij}}{p_{hj}} \right)^{-\rho_j}.
\]
Therefore the Hicksian demands are
\begin{equation}
q_{ij}
=
\xi_i \left( \frac{p_{ij}}{P_j^D} \right)^{-\rho_j} Q_j^D,
\label{eq:app_bottom_demand}
\end{equation}
where
\begin{equation}
P_j^D
=
\left(
\sum_{i=1}^{N_j}
\xi_i p_{ij}^{1-\rho_j}
\right)^{\frac{1}{1-\rho_j}}
\label{eq:app_bottom_price_index}
\end{equation}
is the CES price index for the domestic composite.

\subsubsection{Inverse demand for a domestic variety}

Rearranging \eqref{eq:app_bottom_demand},
\begin{equation}
p_{ij}
=
P_j^D
\left( \frac{q_{ij}}{\xi_i Q_j^D} \right)^{-\frac{1}{\rho_j}}
=
\xi_i^{\frac{1}{\rho_j}} P_j^D \left( \frac{q_{ij}}{Q_j^D} \right)^{-\frac{1}{\rho_j}}.
\label{eq:app_bottom_inverse}
\end{equation}
Taking logs gives
\begin{equation}
\ln p_{ij}
=
\ln P_j^D
- \frac{1}{\rho_j} \ln q_{ij}
+ \frac{1}{\rho_j} \ln Q_j^D
+ \frac{1}{\rho_j} \ln \xi_i.
\label{eq:app_bottom_inverse_logs}
\end{equation}

\subsubsection{Revenue shares}

Multiply \eqref{eq:app_bottom_demand} by $p_{ij}$ and divide by total domestic expenditure $P_j^D Q_j^D$:
\begin{equation}
\tilde{s}_{ij}
\equiv
\frac{p_{ij} q_{ij}}{P_j^D Q_j^D}
=
\xi_i \left( \frac{p_{ij}}{P_j^D} \right)^{1-\rho_j}.
\label{eq:app_within_share}
\end{equation}
This is firm $i$'s expenditure share within the \emph{domestic} segment of sector $j$.

Its share in \emph{total} sectoral expenditure is
\begin{equation}
s_{ij}
\equiv
\frac{p_{ij} q_{ij}}{P_j Q_j}
=
\frac{P_j^D Q_j^D}{P_j Q_j}
\cdot
\frac{p_{ij} q_{ij}}{P_j^D Q_j^D}
=
\lambda_j^D \tilde{s}_{ij}.
\label{eq:app_total_share}
\end{equation}
If one uses the reduced-form notation of the main text that absorbs relative-price terms into the domestic share, \eqref{eq:app_total_share} becomes $s_{ij} = \Lambda_j^D \tilde{s}_{ij}$.

\subsection{Useful derivative identities}

The next step is to compute how an increase in firm $i$'s own output $q_{ij}$ affects the relevant aggregates $Q_j^D$ and $Q_j$. These derivatives are what allow the three CES nests to show up inside the perceived elasticity.

\subsubsection{How a change in $q_{ij}$ affects the domestic composite $Q_j^D$}

Differentiate the domestic aggregator with respect to $q_{ij}$:
\[
\frac{\partial Q_j^D}{\partial q_{ij}}
=
(Q_j^D)^{\frac{1}{\rho_j}}
\xi_i^{\frac{1}{\rho_j}}
q_{ij}^{-\frac{1}{\rho_j}}.
\]
Multiply by $q_{ij}/Q_j^D$:
\[
\frac{\partial \ln Q_j^D}{\partial \ln q_{ij}}
=
\frac{q_{ij}}{Q_j^D}
\frac{\partial Q_j^D}{\partial q_{ij}}
=
\xi_i^{\frac{1}{\rho_j}} q_{ij}^{\frac{\rho_j-1}{\rho_j}} (Q_j^D)^{-\frac{\rho_j-1}{\rho_j}}.
\]
Now use the inverse-demand equation \eqref{eq:app_bottom_inverse}:
\[
\frac{p_{ij} q_{ij}}{P_j^D Q_j^D}
=
\xi_i^{\frac{1}{\rho_j}} q_{ij}^{\frac{\rho_j-1}{\rho_j}} (Q_j^D)^{-\frac{\rho_j-1}{\rho_j}}
=
\tilde{s}_{ij}.
\]
Therefore
\begin{equation}
\frac{\partial \ln Q_j^D}{\partial \ln q_{ij}} = \tilde{s}_{ij}.
\label{eq:app_dlogQD_dlogq}
\end{equation}
Intuition: if firm $i$ is large within the domestic segment, increasing its own output has a larger effect on the domestic composite.

\subsubsection{How a change in $Q_j^D$ affects the sectoral composite $Q_j$}

Differentiate the middle CES aggregator with respect to $Q_j^D$:
\[
\frac{\partial Q_j}{\partial Q_j^D}
=
Q_j^{\frac{1}{\nu_j}}
(\Lambda_j^D)^{\frac{1}{\nu_j}}
(Q_j^D)^{-\frac{1}{\nu_j}}.
\]
Multiply by $Q_j^D/Q_j$:
\[
\frac{\partial \ln Q_j}{\partial \ln Q_j^D}
=
(\Lambda_j^D)^{\frac{1}{\nu_j}} (Q_j^D)^{\frac{\nu_j-1}{\nu_j}} Q_j^{-\frac{\nu_j-1}{\nu_j}}.
\]
Using \eqref{eq:app_middle_inverse_D}, this expression is exactly the realized domestic expenditure share:
\begin{equation}
\frac{\partial \ln Q_j}{\partial \ln Q_j^D}
=
\frac{P_j^D Q_j^D}{P_j Q_j}
=
\lambda_j^D.
\label{eq:app_dlogQ_dlogQD}
\end{equation}

Combine \eqref{eq:app_dlogQD_dlogq} and \eqref{eq:app_dlogQ_dlogQD}:
\begin{equation}
\frac{\partial \ln Q_j}{\partial \ln q_{ij}}
=
\frac{\partial \ln Q_j}{\partial \ln Q_j^D}
\cdot
\frac{\partial \ln Q_j^D}{\partial \ln q_{ij}}
=
\lambda_j^D \tilde{s}_{ij}
=
s_{ij}.
\label{eq:app_dlogQ_dlogq}
\end{equation}
So the effect of firm $i$'s quantity on the whole sectoral composite is exactly its total sectoral expenditure share.

\subsection{Cost minimization and marginal cost}

Domestic firm $i$ in sector $j$ produces with
\[
q_{ij} = \varphi_{ij} L_{ij}^{\beta_j} M_{ij}^{1-\beta_j},
\qquad 0 < \beta_j < 1.
\]
Given factor prices $w$ and $P_j^M$, the cost-minimization problem is
\[
\min_{L_{ij},M_{ij}}
w L_{ij} + P_j^M M_{ij}
\qquad \text{s.t.} \qquad
q_{ij} = \varphi_{ij} L_{ij}^{\beta_j} M_{ij}^{1-\beta_j}.
\]
The Lagrangian is
\[
\mathcal{L}
=
w L_{ij} + P_j^M M_{ij}
+
\lambda_{ij}
\left[
q_{ij} - \varphi_{ij} L_{ij}^{\beta_j} M_{ij}^{1-\beta_j}
\right].
\]

The FOCs are
\begin{equation}
w
=
\lambda_{ij} \varphi_{ij} \beta_j L_{ij}^{\beta_j-1} M_{ij}^{1-\beta_j},
\label{eq:app_cost_foc_L}
\end{equation}
\begin{equation}
P_j^M
=
\lambda_{ij} \varphi_{ij} (1-\beta_j) L_{ij}^{\beta_j} M_{ij}^{-\beta_j}.
\label{eq:app_cost_foc_M}
\end{equation}

Divide \eqref{eq:app_cost_foc_L} by \eqref{eq:app_cost_foc_M}:
\[
\frac{w}{P_j^M}
=
\frac{\beta_j}{1-\beta_j}
\frac{M_{ij}}{L_{ij}}.
\]
Hence the cost-minimizing input ratio is
\begin{equation}
\frac{M_{ij}}{L_{ij}}
=
\frac{1-\beta_j}{\beta_j} \frac{w}{P_j^M}.
\label{eq:app_input_ratio}
\end{equation}

Now multiply \eqref{eq:app_cost_foc_L} by $L_{ij}$ and \eqref{eq:app_cost_foc_M} by $M_{ij}$:
\[
w L_{ij}
=
\lambda_{ij} \beta_j q_{ij},
\qquad
P_j^M M_{ij}
=
\lambda_{ij} (1-\beta_j) q_{ij}.
\]
Adding the two equations gives
\[
w L_{ij} + P_j^M M_{ij} = \lambda_{ij} q_{ij}.
\]
Because production is constant returns to scale, total cost is proportional to output, so the multiplier $\lambda_{ij}$ is the nominal marginal cost:
\[
\lambda_{ij} = \text{mc}_{ij}.
\]

Solving for the exact unit-cost function yields
\begin{equation}
\text{mc}_{ij}
=
\frac{1}{\varphi_{ij}}
\left( \frac{w}{\beta_j} \right)^{\beta_j}
\left( \frac{P_j^M}{1-\beta_j} \right)^{1-\beta_j}.
\label{eq:app_exact_mc}
\end{equation}
Equivalently,
\[
\text{mc}_{ij}
=
\frac{w^{\beta_j} (P_j^M)^{1-\beta_j}}
{\varphi_{ij} \beta_j^{\beta_j} (1-\beta_j)^{1-\beta_j}}.
\]
Ignoring the sector-specific constant $\beta_j^{-\beta_j} (1-\beta_j)^{-(1-\beta_j)}$, the simpler proportionality used in the main text is
\begin{equation}
\text{mc}_{ij} \propto \varphi_{ij}^{-1} w^{\beta_j} (P_j^M)^{1-\beta_j}.
\label{eq:app_mc_level}
\end{equation}

Taking logs of \eqref{eq:app_exact_mc} or \eqref{eq:app_mc_level} gives
\begin{equation}
d \ln \text{mc}_{ij}
=
- d \ln \varphi_{ij}
+ \beta_j d \ln w
+ (1-\beta_j) d \ln P_j^M.
\label{eq:app_mc_growth}
\end{equation}

\subsection{Input-output propagation of Chinese input shocks}

Let $\mathbf{p}^M$ be the vector of sectoral log input prices, and let $A$ denote the domestic input-coefficient matrix. A first-order approximation of the network structure is
\[
\mathbf{p}^M = A \mathbf{p}^M + \Omega^{CHN} \mathbf{p}^{CHN},
\]
where $\mathbf{p}^{CHN}$ is the vector of Chinese input prices and $\Omega^{CHN}$ collects Chinese sourcing weights.

Rearranging,
\[
(I-A)\mathbf{p}^M = \Omega^{CHN} \mathbf{p}^{CHN},
\qquad
\mathbf{p}^M = (I-A)^{-1} \Omega^{CHN} \mathbf{p}^{CHN}.
\]
Let $L \equiv (I-A)^{-1}$ be the Leontief inverse. For sector $j$,
\begin{equation}
d \ln P_j^M
\approx
\sum_k l_{jk} \omega_k^{CHN} d \ln P_k^{CHN},
\label{eq:app_network_pass}
\end{equation}
where $l_{jk}$ is the $(j,k)$ entry of $L$ and $\omega_k^{CHN}$ is the relevant Chinese sourcing weight.

Substituting \eqref{eq:app_network_pass} into \eqref{eq:app_mc_growth} gives
\begin{equation}
d \ln \text{mc}_{ij}
\approx
- d \ln \varphi_{ij}
+ \beta_j d \ln w
+ (1-\beta_j) \sum_k l_{jk} \omega_k^{CHN} d \ln P_k^{CHN}.
\label{eq:app_mc_reduced}
\end{equation}

\subsection{Deriving the perceived elasticity step by step}

This is the core of the model. Since firms compete in quantities, the cleanest route is:
\begin{enumerate}
    \item write the inverse demand of firm $i$;
    \item differentiate it with respect to $q_{ij}$;
    \item track how a change in $q_{ij}$ moves $Q_j^D$, then $Q_j$, then $P_j$.
\end{enumerate}

\subsubsection{Step 1: inverse demand at the bottom nest}

From \eqref{eq:app_bottom_inverse_logs},
\[
\ln p_{ij}
=
\ln P_j^D
- \frac{1}{\rho_j} \ln q_{ij}
+ \frac{1}{\rho_j} \ln Q_j^D
+ \frac{1}{\rho_j} \ln \xi_i.
\]
Differentiate with respect to $\ln q_{ij}$:
\begin{equation}
\frac{d \ln p_{ij}}{d \ln q_{ij}}
=
\frac{d \ln P_j^D}{d \ln q_{ij}}
- \frac{1}{\rho_j}
+ \frac{1}{\rho_j} \frac{d \ln Q_j^D}{d \ln q_{ij}}.
\label{eq:app_step1}
\end{equation}
Using \eqref{eq:app_dlogQD_dlogq},
\begin{equation}
\frac{d \ln p_{ij}}{d \ln q_{ij}}
=
\frac{d \ln P_j^D}{d \ln q_{ij}}
- \frac{1-\tilde{s}_{ij}}{\rho_j}.
\label{eq:app_step1_simplified}
\end{equation}

\subsubsection{Step 2: how $P_j^D$ responds to $q_{ij}$}

From \eqref{eq:app_middle_inverse_diff},
\[
d \ln P_j^D
=
d \ln P_j
- \frac{1}{\nu_j} d \ln Q_j^D
+ \frac{1}{\nu_j} d \ln Q_j.
\]
Divide by $d \ln q_{ij}$:
\begin{equation}
\frac{d \ln P_j^D}{d \ln q_{ij}}
=
\frac{d \ln P_j}{d \ln q_{ij}}
- \frac{1}{\nu_j} \frac{d \ln Q_j^D}{d \ln q_{ij}}
+ \frac{1}{\nu_j} \frac{d \ln Q_j}{d \ln q_{ij}}.
\label{eq:app_step2}
\end{equation}
Using \eqref{eq:app_dlogQD_dlogq} and \eqref{eq:app_dlogQ_dlogq},
\begin{equation}
\frac{d \ln P_j^D}{d \ln q_{ij}}
=
\frac{d \ln P_j}{d \ln q_{ij}}
- \frac{\tilde{s}_{ij}}{\nu_j}
+ \frac{s_{ij}}{\nu_j}
=
\frac{d \ln P_j}{d \ln q_{ij}}
- \frac{\tilde{s}_{ij} - s_{ij}}{\nu_j}.
\label{eq:app_step2_simplified}
\end{equation}

\subsubsection{Step 3: how $P_j$ responds to $q_{ij}$}

From the sector-level inverse demand \eqref{eq:app_top_inverse},
\[
\ln P_j
=
\ln P
+ \frac{1}{\eta} \ln \alpha_j
- \frac{1}{\eta} \ln Q_j
+ \frac{1}{\eta} \ln C.
\]
Holding the aggregate objects $P$ and $C$ fixed from the viewpoint of an individual firm inside sector $j$,
\begin{equation}
\frac{d \ln P_j}{d \ln q_{ij}}
=
- \frac{1}{\eta} \frac{d \ln Q_j}{d \ln q_{ij}}
=
- \frac{s_{ij}}{\eta},
\label{eq:app_step3}
\end{equation}
where the last equality uses \eqref{eq:app_dlogQ_dlogq}.

\subsubsection{Step 4: combine the three pieces}

Insert \eqref{eq:app_step3} into \eqref{eq:app_step2_simplified}:
\[
\frac{d \ln P_j^D}{d \ln q_{ij}}
=
- \frac{\tilde{s}_{ij} - s_{ij}}{\nu_j}
- \frac{s_{ij}}{\eta}.
\]
Now plug this into \eqref{eq:app_step1_simplified}:
\[
\frac{d \ln p_{ij}}{d \ln q_{ij}}
=
- \frac{1-\tilde{s}_{ij}}{\rho_j}
- \frac{\tilde{s}_{ij} - s_{ij}}{\nu_j}
- \frac{s_{ij}}{\eta}.
\]
Therefore
\begin{equation}
- \frac{d \ln p_{ij}}{d \ln q_{ij}}
=
\frac{1-\tilde{s}_{ij}}{\rho_j}
+
\frac{\tilde{s}_{ij} - s_{ij}}{\nu_j}
+
\frac{s_{ij}}{\eta}.
\label{eq:app_inverse_demand_elasticity}
\end{equation}

Define the firm's \emph{perceived elasticity} under Cournot competition as
\begin{equation}
\frac{1}{\sigma_{ij}^C}
\equiv
- \frac{d \ln p_{ij}}{d \ln q_{ij}}.
\label{eq:app_sigma_definition}
\end{equation}
Equivalently, since price and quantity are locally inverse functions,
\[
\sigma_{ij}^C
=
- \frac{d \ln q_{ij}}{d \ln p_{ij}}.
\]
Combining \eqref{eq:app_inverse_demand_elasticity} and \eqref{eq:app_sigma_definition} yields
\begin{equation}
\sigma_{ij}^C
=
\left(
\frac{1-\tilde{s}_{ij}}{\rho_j}
+
\frac{\tilde{s}_{ij} - s_{ij}}{\nu_j}
+
\frac{s_{ij}}{\eta}
\right)^{-1}.
\label{eq:app_sigma_nested}
\end{equation}

\paragraph{Interpretation.}
Equation \eqref{eq:app_sigma_nested} is easiest to read term by term.
\begin{itemize}
    \item The term $\frac{1-\tilde{s}_{ij}}{\rho_j}$ comes from substitution across \emph{domestic varieties}. If the firm is small within the domestic segment, this force is strong.
    \item The term $\frac{\tilde{s}_{ij}-s_{ij}}{\nu_j}$ comes from substitution between the \emph{domestic composite} and the \emph{foreign composite} inside the sector.
    \item The term $\frac{s_{ij}}{\eta}$ comes from substitution across \emph{sectors} in final demand.
\end{itemize}
Because $\rho_j > \nu_j > \eta$, the inverse elasticities satisfy
\[
\frac{1}{\rho_j} < \frac{1}{\nu_j} < \frac{1}{\eta}.
\]
So as a firm becomes more important in broader and broader aggregates, the relevant residual demand it faces becomes less elastic.

\subsection{Cournot firm's FOC and the markup formula}

Firm $i$ chooses quantity $q_{ij}$ to maximize
\[
\pi_{ij}
=
\bigl(p_{ij}(q_{ij},q_{-ij}) - \text{mc}_{ij}\bigr) q_{ij} - f_{ij}.
\]
The first-order condition is
\begin{equation}
\frac{\partial \pi_{ij}}{\partial q_{ij}}
=
p_{ij}
+ q_{ij} \frac{\partial p_{ij}}{\partial q_{ij}}
- \text{mc}_{ij}
=
0.
\label{eq:app_cournot_foc}
\end{equation}

Divide \eqref{eq:app_cournot_foc} by $p_{ij}$:
\[
1
+ \frac{q_{ij}}{p_{ij}} \frac{\partial p_{ij}}{\partial q_{ij}}
- \frac{\text{mc}_{ij}}{p_{ij}}
=
0.
\]
Using the definition of the perceived elasticity,
\[
\frac{q_{ij}}{p_{ij}} \frac{\partial p_{ij}}{\partial q_{ij}}
=
\frac{d \ln p_{ij}}{d \ln q_{ij}}
=
- \frac{1}{\sigma_{ij}^C}.
\]
Hence
\[
1 - \frac{1}{\sigma_{ij}^C} = \frac{\text{mc}_{ij}}{p_{ij}}.
\]
Define the markup
\[
\mu_{ij} \equiv \frac{p_{ij}}{\text{mc}_{ij}}.
\]
Then
\begin{equation}
\frac{1}{\mu_{ij}}
=
1 - \frac{1}{\sigma_{ij}^C},
\qquad
\mu_{ij}
=
\frac{1}{1-\frac{1}{\sigma_{ij}^C}}
=
\frac{\sigma_{ij}^C}{\sigma_{ij}^C-1}.
\label{eq:app_markup_sigma}
\end{equation}

Substituting \eqref{eq:app_sigma_nested} into \eqref{eq:app_markup_sigma},
\begin{equation}
\frac{1}{\mu_{ij}}
=
1
-
\left[
\frac{1-\tilde{s}_{ij}}{\rho_j}
+
\frac{\tilde{s}_{ij} - s_{ij}}{\nu_j}
+
\frac{s_{ij}}{\eta}
\right].
\label{eq:app_inverse_markup_unsimplified}
\end{equation}
This is already a useful formula: a larger firm faces a less elastic residual demand, so $1/\mu_{ij}$ is lower and $\mu_{ij}$ is higher.

Now use $s_{ij} = \lambda_j^D \tilde{s}_{ij}$, so $\tilde{s}_{ij} = s_{ij}/\lambda_j^D$. Then \eqref{eq:app_inverse_markup_unsimplified} becomes
\begin{align}
\frac{1}{\mu_{ij}}
&=
1 - \frac{1}{\rho_j}
+ \frac{\tilde{s}_{ij}}{\rho_j}
- \frac{\tilde{s}_{ij}}{\nu_j}
+ \frac{s_{ij}}{\nu_j}
- \frac{s_{ij}}{\eta}
\notag \\
&=
\left( 1 - \frac{1}{\rho_j} \right)
-
\left[
\frac{1}{\lambda_j^D}
\left( \frac{1}{\nu_j} - \frac{1}{\rho_j} \right)
+
\left( \frac{1}{\eta} - \frac{1}{\nu_j} \right)
\right] s_{ij}.
\label{eq:app_inverse_markup_linear}
\end{align}

\paragraph{Important notation point.}
Equation \eqref{eq:app_inverse_markup_linear} is a linear equation for the \emph{inverse markup} $1/\mu_{ij}$, not for $\mu_{ij}$ itself. Therefore the corresponding markup is
\begin{equation}
\mu_{ij}
=
\frac{1}{
\left( 1 - \frac{1}{\rho_j} \right)
-
\left[
\frac{1}{\lambda_j^D}
\left( \frac{1}{\nu_j} - \frac{1}{\rho_j} \right)
+
\left( \frac{1}{\eta} - \frac{1}{\nu_j} \right)
\right] s_{ij}
}.
\label{eq:app_markup_full}
\end{equation}
In the main text, I suppress the distinction between $\lambda_j^D$ and $\Lambda_j^D$ for expositional simplicity; under that notation, \eqref{eq:app_inverse_markup_linear} becomes exactly the linear markup condition reported there, except that it should strictly be read as an inverse-markup condition.

\subsection{Sectoral markup and concentration}

Define the sectoral markup as the share-weighted harmonic mean of firm markups:
\begin{equation}
\mu_j
\equiv
\left(
\sum_{i=1}^{N_j} \frac{s_{ij}}{\mu_{ij}}
\right)^{-1}.
\label{eq:app_sector_harmonic}
\end{equation}
Taking inverses,
\[
\frac{1}{\mu_j}
=
\sum_{i=1}^{N_j} s_{ij} \frac{1}{\mu_{ij}}.
\]
Now substitute \eqref{eq:app_inverse_markup_linear}:
\begin{align}
\frac{1}{\mu_j}
&=
\sum_{i=1}^{N_j}
s_{ij}
\left[
\left( 1 - \frac{1}{\rho_j} \right)
-
\left(
\frac{1}{\lambda_j^D}
\left( \frac{1}{\nu_j} - \frac{1}{\rho_j} \right)
+
\left( \frac{1}{\eta} - \frac{1}{\nu_j} \right)
\right) s_{ij}
\right]
\notag \\
&=
\left( 1 - \frac{1}{\rho_j} \right) \sum_{i=1}^{N_j} s_{ij}
-
\left[
\frac{1}{\lambda_j^D}
\left( \frac{1}{\nu_j} - \frac{1}{\rho_j} \right)
+
\left( \frac{1}{\eta} - \frac{1}{\nu_j} \right)
\right]
\sum_{i=1}^{N_j} s_{ij}^2.
\label{eq:app_sector_inverse_step}
\end{align}
Since $\sum_i s_{ij} = 1$, and defining the Herfindahl index
\[
\text{HHI}_j \equiv \sum_{i=1}^{N_j} s_{ij}^2,
\]
we obtain
\begin{equation}
\frac{1}{\mu_j}
=
\left( 1 - \frac{1}{\rho_j} \right)
-
\left[
\frac{1}{\lambda_j^D}
\left( \frac{1}{\nu_j} - \frac{1}{\rho_j} \right)
+
\left( \frac{1}{\eta} - \frac{1}{\nu_j} \right)
\right]
\text{HHI}_j.
\label{eq:app_sector_inverse_markup}
\end{equation}
Therefore
\begin{equation}
\mu_j
=
\frac{1}{
\left( 1 - \frac{1}{\rho_j} \right)
-
\left[
\frac{1}{\lambda_j^D}
\left( \frac{1}{\nu_j} - \frac{1}{\rho_j} \right)
+
\left( \frac{1}{\eta} - \frac{1}{\nu_j} \right)
\right]
\text{HHI}_j
}.
\label{eq:app_sector_markup}
\end{equation}

Again, the \emph{linear} object is $1/\mu_j$, not $\mu_j$.

\subsubsection{Why concentration matters}

Let
\[
A_j \equiv 1 - \frac{1}{\rho_j},
\qquad
B_j \equiv
\frac{1}{\lambda_j^D}
\left( \frac{1}{\nu_j} - \frac{1}{\rho_j} \right)
+
\left( \frac{1}{\eta} - \frac{1}{\nu_j} \right).
\]
Then
\[
\frac{1}{\mu_j} = A_j - B_j \text{HHI}_j,
\qquad
\mu_j = \frac{1}{A_j - B_j \text{HHI}_j}.
\]
Differentiating,
\[
\frac{\partial \mu_j}{\partial \text{HHI}_j}
=
\frac{B_j}{(A_j - B_j \text{HHI}_j)^2}.
\]
Under $\rho_j > \nu_j > \eta$, we have
\[
\frac{1}{\nu_j} - \frac{1}{\rho_j} > 0,
\qquad
\frac{1}{\eta} - \frac{1}{\nu_j} > 0,
\]
so $B_j > 0$. Hence
\begin{equation}
\frac{\partial \mu_j}{\partial \text{HHI}_j} > 0.
\label{eq:app_markup_concentration_positive}
\end{equation}
Sectoral markups rise with concentration because more expenditure gets concentrated on firms that face less elastic residual demand.

\subsection{Dynamic extension and endogenous exit}

Let $x_{ijt}$ denote the state of firm $i$ in sector $j$ at date $t$. This state includes productivity, demand shifters, lagged market share, and the output- and input-trade states $(\tau_{jt}^O,\tau_{jt}^I)$.

If the firm stays active, its value function is
\begin{equation}
V(x_{ijt})
=
\max
\left\{
0,
\pi(x_{ijt})
+ \beta \mathbb{E}\left[ V(x_{ij,t+1}) \mid x_{ijt} \right]
\right\},
\label{eq:app_bellman}
\end{equation}
where $\beta \in (0,1)$ is the discount factor.

Define the continuation surplus
\[
W(x_{ijt})
\equiv
\pi(x_{ijt})
+ \beta \mathbb{E}\left[ V(x_{ij,t+1}) \mid x_{ijt} \right].
\]
The optimal exit rule is
\begin{equation}
\text{Exit}_{ijt} = 1
\quad \Longleftrightarrow \quad
W(x_{ijt}) < 0.
\label{eq:app_exit_rule}
\end{equation}
Equivalently, if current profitability can be summarized by a latent scalar $\theta_{ijt}$, there exists a cutoff function $\theta^*(x_{ijt})$ such that
\begin{equation}
\text{Exit}_{ijt} = 1
\quad \Longleftrightarrow \quad
\theta_{ijt} < \theta^*(x_{ijt}).
\label{eq:app_exit_threshold}
\end{equation}

\paragraph{Why this matters empirically.}
Trade shocks affect both current profits and future continuation values. A rise in final-good import competition may reduce current operating profits; a favorable input-supply shock may reduce marginal cost and increase the value of staying active. Therefore the sample of surviving firms is not random. Observed firm-level markup changes mix:
\begin{itemize}
    \item genuine within-incumbent adjustment;
    \item reallocation across surviving incumbents;
    \item selection induced by endogenous exit.
\end{itemize}
This is the reason the empirical section later introduces a selection-correction strategy.
.

\section{Theoretical Appendix}
\label{app:theory}

This appendix solves the model in Section 2 step by step. I write the algebra in a way that should be readable for a good undergraduate student who is already comfortable with constrained optimization, CES demand systems, and the idea of a markup as price over marginal cost.

The main goals are:
\begin{enumerate}
    \item derive the demand system at each CES nest from first principles;
    \item derive the inverse-demand system because firms compete in quantities;
    \item show carefully how the perceived elasticity is built from the three substitution margins in the model;
    \item derive the Cournot markup formula and the sectoral markup aggregator;
    \item clarify one notation point: the linear expression in market shares is the \emph{inverse markup} $1/\mu_{ij}$, not the markup $\mu_{ij}$ itself.
\end{enumerate}

\subsection{A useful CES template}

Before solving each layer separately, it is useful to record one generic CES result that we will reuse.

Suppose a composite $X$ is built from a collection of goods $\{x_n\}_{n \in \mathcal{N}}$:
\[
X =
\left(
\sum_{n \in \mathcal{N}}
a_n^{\frac{1}{\sigma}} x_n^{\frac{\sigma-1}{\sigma}}
\right)^{\frac{\sigma}{\sigma-1}},
\qquad \sigma > 1.
\]
Given prices $\{p_n\}$, the dual expenditure-minimization problem is
\[
\min_{\{x_n\}} \sum_{n \in \mathcal{N}} p_n x_n
\qquad \text{s.t.} \qquad
X \geq \bar{X}.
\]
The Lagrangian is
\[
\mathcal{L}
=
\sum_{n \in \mathcal{N}} p_n x_n
+
\lambda
\left[
\bar{X}
-
\left(
\sum_{n \in \mathcal{N}}
a_n^{\frac{1}{\sigma}} x_n^{\frac{\sigma-1}{\sigma}}
\right)^{\frac{\sigma}{\sigma-1}}
\right].
\]

Let
\[
M \equiv \sum_{n \in \mathcal{N}} a_n^{\frac{1}{\sigma}} x_n^{\frac{\sigma-1}{\sigma}},
\qquad
X = M^{\frac{\sigma}{\sigma-1}}.
\]
Then
\[
\frac{\partial X}{\partial x_n}
=
\frac{\sigma}{\sigma-1}
M^{\frac{1}{\sigma-1}}
\cdot
a_n^{\frac{1}{\sigma}}
\cdot
\frac{\sigma-1}{\sigma} x_n^{-\frac{1}{\sigma}}
=
X^{\frac{1}{\sigma}} a_n^{\frac{1}{\sigma}} x_n^{-\frac{1}{\sigma}}.
\]
Therefore the first-order condition for each $n$ is
\begin{equation}
p_n = \lambda X^{\frac{1}{\sigma}} a_n^{\frac{1}{\sigma}} x_n^{-\frac{1}{\sigma}}.
\label{eq:app_generic_ces_foc}
\end{equation}

Take two goods $n$ and $m$. Dividing the two FOCs gives
\[
\frac{p_n}{p_m}
=
\left( \frac{a_n}{a_m} \right)^{\frac{1}{\sigma}}
\left( \frac{x_n}{x_m} \right)^{-\frac{1}{\sigma}}
\Longrightarrow
\frac{x_n}{x_m}
=
\frac{a_n}{a_m}
\left( \frac{p_n}{p_m} \right)^{-\sigma}.
\]
This is the familiar CES relative-demand equation. Solving completely gives
\begin{equation}
x_n = a_n \left( \frac{p_n}{P_X} \right)^{-\sigma} X,
\label{eq:app_generic_hicksian}
\end{equation}
where the CES price index is
\begin{equation}
P_X =
\left(
\sum_{n \in \mathcal{N}}
a_n p_n^{1-\sigma}
\right)^{\frac{1}{1-\sigma}}.
\label{eq:app_generic_price_index}
\end{equation}

Finally, rearranging \eqref{eq:app_generic_hicksian} gives the inverse-demand form
\begin{equation}
p_n
=
P_X
\left( \frac{x_n}{a_n X} \right)^{-\frac{1}{\sigma}}
=
a_n^{\frac{1}{\sigma}} P_X \left( \frac{x_n}{X} \right)^{-\frac{1}{\sigma}}.
\label{eq:app_generic_inverse_demand}
\end{equation}

We now apply the same steps to each nest in the model.

\subsection{Top nest: allocation across sectors}

The representative household values sectoral composites according to
\[
C =
\left(
\int_0^1
\alpha_j^{\frac{1}{\eta}} Q_j^{\frac{\eta-1}{\eta}} dj
\right)^{\frac{\eta}{\eta-1}},
\qquad
\eta > 1,
\qquad
\int_0^1 \alpha_j dj = 1.
\]

\subsubsection{Expenditure minimization problem}

Fix a target final consumption level $\bar{C}$. The household solves
\[
\min_{\{Q_j\}_{j \in [0,1]}}
\int_0^1 P_j Q_j dj
\qquad \text{s.t.} \qquad
\left(
\int_0^1
\alpha_j^{\frac{1}{\eta}} Q_j^{\frac{\eta-1}{\eta}} dj
\right)^{\frac{\eta}{\eta-1}}
\geq \bar{C}.
\]

The Lagrangian is
\[
\mathcal{L}
=
\int_0^1 P_j Q_j dj
+
\lambda
\left[
\bar{C}
-
\left(
\int_0^1
\alpha_j^{\frac{1}{\eta}} Q_j^{\frac{\eta-1}{\eta}} dj
\right)^{\frac{\eta}{\eta-1}}
\right].
\]

Using the generic CES template with a continuum of goods, the FOC for sector $j$ is
\begin{equation}
P_j = \lambda C^{\frac{1}{\eta}} \alpha_j^{\frac{1}{\eta}} Q_j^{-\frac{1}{\eta}}.
\label{eq:app_top_foc}
\end{equation}
Take two sectors $j$ and $\ell$. Dividing their FOCs gives
\[
\frac{Q_j}{Q_\ell}
=
\frac{\alpha_j}{\alpha_\ell}
\left( \frac{P_j}{P_\ell} \right)^{-\eta}.
\]
Hence sectoral demand is
\begin{equation}
Q_j = \alpha_j \left( \frac{P_j}{P} \right)^{-\eta} C,
\label{eq:app_top_demand}
\end{equation}
where
\begin{equation}
P =
\left(
\int_0^1 \alpha_j P_j^{1-\eta} dj
\right)^{\frac{1}{1-\eta}}
\label{eq:app_top_price_index}
\end{equation}
is the final CES price index.

\subsubsection{Sectoral expenditure}

Multiply \eqref{eq:app_top_demand} by $P_j$:
\[
D_j \equiv P_j Q_j = \alpha_j \left( \frac{P_j}{P} \right)^{1-\eta} P C.
\]
Using $Y = PC$, we obtain
\begin{equation}
D_j = \alpha_j \left( \frac{P_j}{P} \right)^{1-\eta} Y.
\label{eq:app_sector_expenditure}
\end{equation}
This is the sector-level expenditure expression reported in the main text.

\subsubsection{Inverse demand at the sector level}

Rearranging \eqref{eq:app_top_demand} gives
\begin{equation}
P_j = P \, \alpha_j^{\frac{1}{\eta}} \left( \frac{Q_j}{C} \right)^{-\frac{1}{\eta}}.
\label{eq:app_top_inverse}
\end{equation}
Taking logs,
\[
\ln P_j = \ln P + \frac{1}{\eta} \ln \alpha_j - \frac{1}{\eta} \ln Q_j + \frac{1}{\eta} \ln C.
\]
So, holding $P$, $\alpha_j$, and $C$ fixed,
\begin{equation}
\frac{d \ln P_j}{d \ln Q_j} = - \frac{1}{\eta}.
\label{eq:app_top_inverse_elasticity}
\end{equation}
This is the inverse-demand elasticity at the top nest.

\subsection{Middle nest: domestic versus foreign supply inside sector}

Within sector $j$, the sectoral composite is
\[
Q_j =
\left[
(\Lambda_j^D)^{\frac{1}{\nu_j}} (Q_j^D)^{\frac{\nu_j-1}{\nu_j}}
+
(1-\Lambda_j^D)^{\frac{1}{\nu_j}} (Q_j^F)^{\frac{\nu_j-1}{\nu_j}}
\right]^{\frac{\nu_j}{\nu_j-1}},
\qquad \nu_j > 1.
\]

\subsubsection{Expenditure minimization problem}

For a given target $\bar{Q}_j$, minimize expenditure on the domestic and foreign composites:
\[
\min_{Q_j^D,Q_j^F}
P_j^D Q_j^D + P_j^F Q_j^F
\qquad \text{s.t.} \qquad
Q_j \geq \bar{Q}_j.
\]
The Lagrangian is
\[
\mathcal{L}
=
P_j^D Q_j^D + P_j^F Q_j^F
+
\lambda_j
\left[
\bar{Q}_j
-
\left[
(\Lambda_j^D)^{\frac{1}{\nu_j}} (Q_j^D)^{\frac{\nu_j-1}{\nu_j}}
+
(1-\Lambda_j^D)^{\frac{1}{\nu_j}} (Q_j^F)^{\frac{\nu_j-1}{\nu_j}}
\right]^{\frac{\nu_j}{\nu_j-1}}
\right].
\]

The FOCs are
\begin{equation}
P_j^D
=
\lambda_j Q_j^{\frac{1}{\nu_j}}
(\Lambda_j^D)^{\frac{1}{\nu_j}}
(Q_j^D)^{-\frac{1}{\nu_j}},
\label{eq:app_middle_foc_D}
\end{equation}
\begin{equation}
P_j^F
=
\lambda_j Q_j^{\frac{1}{\nu_j}}
(1-\Lambda_j^D)^{\frac{1}{\nu_j}}
(Q_j^F)^{-\frac{1}{\nu_j}}.
\label{eq:app_middle_foc_F}
\end{equation}

Dividing \eqref{eq:app_middle_foc_D} by \eqref{eq:app_middle_foc_F} gives
\[
\frac{Q_j^D}{Q_j^F}
=
\frac{\Lambda_j^D}{1-\Lambda_j^D}
\left( \frac{P_j^D}{P_j^F} \right)^{-\nu_j}.
\]
Therefore the Hicksian demands are
\begin{equation}
Q_j^D = \Lambda_j^D \left( \frac{P_j^D}{P_j} \right)^{-\nu_j} Q_j,
\label{eq:app_middle_demand_D}
\end{equation}
\begin{equation}
Q_j^F = (1-\Lambda_j^D) \left( \frac{P_j^F}{P_j} \right)^{-\nu_j} Q_j,
\label{eq:app_middle_demand_F}
\end{equation}
with price index
\begin{equation}
P_j =
\left[
\Lambda_j^D (P_j^D)^{1-\nu_j}
+
(1-\Lambda_j^D) (P_j^F)^{1-\nu_j}
\right]^{\frac{1}{1-\nu_j}}.
\label{eq:app_middle_price_index}
\end{equation}

\subsubsection{A useful notation point}

The CES weight $\Lambda_j^D$ is not, in general, the same thing as the realized domestic expenditure share. The realized share is
\begin{equation}
\lambda_j^D
\equiv
\frac{P_j^D Q_j^D}{P_j Q_j}
=
\Lambda_j^D \left( \frac{P_j^D}{P_j} \right)^{1-\nu_j}.
\label{eq:app_realized_domestic_share}
\end{equation}
In the main text, I often use $\Lambda_j^D$ as reduced-form shorthand for the domestic expenditure share. For the derivations below, it is cleaner to keep the distinction:
\begin{itemize}
    \item $\Lambda_j^D$: CES weight in preferences;
    \item $\lambda_j^D$: realized domestic expenditure share.
\end{itemize}

\subsubsection{Inverse demand for the domestic composite}

Rearranging \eqref{eq:app_middle_demand_D} gives
\begin{equation}
P_j^D
=
P_j
\left( \frac{Q_j^D}{\Lambda_j^D Q_j} \right)^{-\frac{1}{\nu_j}}.
\label{eq:app_middle_inverse_D}
\end{equation}
Taking logs,
\[
\ln P_j^D
=
\ln P_j
- \frac{1}{\nu_j} \ln Q_j^D
+ \frac{1}{\nu_j} \ln Q_j
+ \frac{1}{\nu_j} \ln \Lambda_j^D.
\]
Hence, when $\Lambda_j^D$ is treated as given,
\begin{equation}
d \ln P_j^D
=
d \ln P_j
- \frac{1}{\nu_j} d \ln Q_j^D
+ \frac{1}{\nu_j} d \ln Q_j.
\label{eq:app_middle_inverse_diff}
\end{equation}

\subsection{Bottom nest: domestic varieties inside the domestic composite}

Within the domestic nest of sector $j$, the domestic composite is
\[
Q_j^D =
\left(
\sum_{i=1}^{N_j}
\xi_i^{\frac{1}{\rho_j}} q_{ij}^{\frac{\rho_j-1}{\rho_j}}
\right)^{\frac{\rho_j}{\rho_j-1}},
\qquad \rho_j > 1.
\]

\subsubsection{Expenditure minimization problem}

For a given target $\bar{Q}_j^D$, minimize spending across domestic varieties:
\[
\min_{\{q_{ij}\}_{i=1}^{N_j}}
\sum_{i=1}^{N_j} p_{ij} q_{ij}
\qquad \text{s.t.} \qquad
Q_j^D \geq \bar{Q}_j^D.
\]
The Lagrangian is
\[
\mathcal{L}
=
\sum_{i=1}^{N_j} p_{ij} q_{ij}
+
\kappa_j
\left[
\bar{Q}_j^D
-
\left(
\sum_{i=1}^{N_j}
\xi_i^{\frac{1}{\rho_j}} q_{ij}^{\frac{\rho_j-1}{\rho_j}}
\right)^{\frac{\rho_j}{\rho_j-1}}
\right].
\]

The FOC for variety $i$ is
\begin{equation}
p_{ij}
=
\kappa_j (Q_j^D)^{\frac{1}{\rho_j}}
\xi_i^{\frac{1}{\rho_j}}
q_{ij}^{-\frac{1}{\rho_j}}.
\label{eq:app_bottom_foc}
\end{equation}
Take two varieties $i$ and $h$ and divide the FOCs:
\[
\frac{q_{ij}}{q_{hj}}
=
\frac{\xi_i}{\xi_h}
\left( \frac{p_{ij}}{p_{hj}} \right)^{-\rho_j}.
\]
Therefore the Hicksian demands are
\begin{equation}
q_{ij}
=
\xi_i \left( \frac{p_{ij}}{P_j^D} \right)^{-\rho_j} Q_j^D,
\label{eq:app_bottom_demand}
\end{equation}
where
\begin{equation}
P_j^D
=
\left(
\sum_{i=1}^{N_j}
\xi_i p_{ij}^{1-\rho_j}
\right)^{\frac{1}{1-\rho_j}}
\label{eq:app_bottom_price_index}
\end{equation}
is the CES price index for the domestic composite.

\subsubsection{Inverse demand for a domestic variety}

Rearranging \eqref{eq:app_bottom_demand},
\begin{equation}
p_{ij}
=
P_j^D
\left( \frac{q_{ij}}{\xi_i Q_j^D} \right)^{-\frac{1}{\rho_j}}
=
\xi_i^{\frac{1}{\rho_j}} P_j^D \left( \frac{q_{ij}}{Q_j^D} \right)^{-\frac{1}{\rho_j}}.
\label{eq:app_bottom_inverse}
\end{equation}
Taking logs gives
\begin{equation}
\ln p_{ij}
=
\ln P_j^D
- \frac{1}{\rho_j} \ln q_{ij}
+ \frac{1}{\rho_j} \ln Q_j^D
+ \frac{1}{\rho_j} \ln \xi_i.
\label{eq:app_bottom_inverse_logs}
\end{equation}

\subsubsection{Revenue shares}

Multiply \eqref{eq:app_bottom_demand} by $p_{ij}$ and divide by total domestic expenditure $P_j^D Q_j^D$:
\begin{equation}
\tilde{s}_{ij}
\equiv
\frac{p_{ij} q_{ij}}{P_j^D Q_j^D}
=
\xi_i \left( \frac{p_{ij}}{P_j^D} \right)^{1-\rho_j}.
\label{eq:app_within_share}
\end{equation}
This is firm $i$'s expenditure share within the \emph{domestic} segment of sector $j$.

Its share in \emph{total} sectoral expenditure is
\begin{equation}
s_{ij}
\equiv
\frac{p_{ij} q_{ij}}{P_j Q_j}
=
\frac{P_j^D Q_j^D}{P_j Q_j}
\cdot
\frac{p_{ij} q_{ij}}{P_j^D Q_j^D}
=
\lambda_j^D \tilde{s}_{ij}.
\label{eq:app_total_share}
\end{equation}
If one uses the reduced-form notation of the main text that absorbs relative-price terms into the domestic share, \eqref{eq:app_total_share} becomes $s_{ij} = \Lambda_j^D \tilde{s}_{ij}$.

\subsection{Useful derivative identities}

The next step is to compute how an increase in firm $i$'s own output $q_{ij}$ affects the relevant aggregates $Q_j^D$ and $Q_j$. These derivatives are what allow the three CES nests to show up inside the perceived elasticity.

\subsubsection{How a change in $q_{ij}$ affects the domestic composite $Q_j^D$}

Differentiate the domestic aggregator with respect to $q_{ij}$:
\[
\frac{\partial Q_j^D}{\partial q_{ij}}
=
(Q_j^D)^{\frac{1}{\rho_j}}
\xi_i^{\frac{1}{\rho_j}}
q_{ij}^{-\frac{1}{\rho_j}}.
\]
Multiply by $q_{ij}/Q_j^D$:
\[
\frac{\partial \ln Q_j^D}{\partial \ln q_{ij}}
=
\frac{q_{ij}}{Q_j^D}
\frac{\partial Q_j^D}{\partial q_{ij}}
=
\xi_i^{\frac{1}{\rho_j}} q_{ij}^{\frac{\rho_j-1}{\rho_j}} (Q_j^D)^{-\frac{\rho_j-1}{\rho_j}}.
\]
Now use the inverse-demand equation \eqref{eq:app_bottom_inverse}:
\[
\frac{p_{ij} q_{ij}}{P_j^D Q_j^D}
=
\xi_i^{\frac{1}{\rho_j}} q_{ij}^{\frac{\rho_j-1}{\rho_j}} (Q_j^D)^{-\frac{\rho_j-1}{\rho_j}}
=
\tilde{s}_{ij}.
\]
Therefore
\begin{equation}
\frac{\partial \ln Q_j^D}{\partial \ln q_{ij}} = \tilde{s}_{ij}.
\label{eq:app_dlogQD_dlogq}
\end{equation}
Intuition: if firm $i$ is large within the domestic segment, increasing its own output has a larger effect on the domestic composite.

\subsubsection{How a change in $Q_j^D$ affects the sectoral composite $Q_j$}

Differentiate the middle CES aggregator with respect to $Q_j^D$:
\[
\frac{\partial Q_j}{\partial Q_j^D}
=
Q_j^{\frac{1}{\nu_j}}
(\Lambda_j^D)^{\frac{1}{\nu_j}}
(Q_j^D)^{-\frac{1}{\nu_j}}.
\]
Multiply by $Q_j^D/Q_j$:
\[
\frac{\partial \ln Q_j}{\partial \ln Q_j^D}
=
(\Lambda_j^D)^{\frac{1}{\nu_j}} (Q_j^D)^{\frac{\nu_j-1}{\nu_j}} Q_j^{-\frac{\nu_j-1}{\nu_j}}.
\]
Using \eqref{eq:app_middle_inverse_D}, this expression is exactly the realized domestic expenditure share:
\begin{equation}
\frac{\partial \ln Q_j}{\partial \ln Q_j^D}
=
\frac{P_j^D Q_j^D}{P_j Q_j}
=
\lambda_j^D.
\label{eq:app_dlogQ_dlogQD}
\end{equation}

Combine \eqref{eq:app_dlogQD_dlogq} and \eqref{eq:app_dlogQ_dlogQD}:
\begin{equation}
\frac{\partial \ln Q_j}{\partial \ln q_{ij}}
=
\frac{\partial \ln Q_j}{\partial \ln Q_j^D}
\cdot
\frac{\partial \ln Q_j^D}{\partial \ln q_{ij}}
=
\lambda_j^D \tilde{s}_{ij}
=
s_{ij}.
\label{eq:app_dlogQ_dlogq}
\end{equation}
So the effect of firm $i$'s quantity on the whole sectoral composite is exactly its total sectoral expenditure share.

\subsection{Cost minimization and marginal cost}

Domestic firm $i$ in sector $j$ produces with
\[
q_{ij} = \varphi_{ij} L_{ij}^{\beta_j} M_{ij}^{1-\beta_j},
\qquad 0 < \beta_j < 1.
\]
Given factor prices $w$ and $P_j^M$, the cost-minimization problem is
\[
\min_{L_{ij},M_{ij}}
w L_{ij} + P_j^M M_{ij}
\qquad \text{s.t.} \qquad
q_{ij} = \varphi_{ij} L_{ij}^{\beta_j} M_{ij}^{1-\beta_j}.
\]
The Lagrangian is
\[
\mathcal{L}
=
w L_{ij} + P_j^M M_{ij}
+
\lambda_{ij}
\left[
q_{ij} - \varphi_{ij} L_{ij}^{\beta_j} M_{ij}^{1-\beta_j}
\right].
\]

The FOCs are
\begin{equation}
w
=
\lambda_{ij} \varphi_{ij} \beta_j L_{ij}^{\beta_j-1} M_{ij}^{1-\beta_j},
\label{eq:app_cost_foc_L}
\end{equation}
\begin{equation}
P_j^M
=
\lambda_{ij} \varphi_{ij} (1-\beta_j) L_{ij}^{\beta_j} M_{ij}^{-\beta_j}.
\label{eq:app_cost_foc_M}
\end{equation}

Divide \eqref{eq:app_cost_foc_L} by \eqref{eq:app_cost_foc_M}:
\[
\frac{w}{P_j^M}
=
\frac{\beta_j}{1-\beta_j}
\frac{M_{ij}}{L_{ij}}.
\]
Hence the cost-minimizing input ratio is
\begin{equation}
\frac{M_{ij}}{L_{ij}}
=
\frac{1-\beta_j}{\beta_j} \frac{w}{P_j^M}.
\label{eq:app_input_ratio}
\end{equation}

Now multiply \eqref{eq:app_cost_foc_L} by $L_{ij}$ and \eqref{eq:app_cost_foc_M} by $M_{ij}$:
\[
w L_{ij}
=
\lambda_{ij} \beta_j q_{ij},
\qquad
P_j^M M_{ij}
=
\lambda_{ij} (1-\beta_j) q_{ij}.
\]
Adding the two equations gives
\[
w L_{ij} + P_j^M M_{ij} = \lambda_{ij} q_{ij}.
\]
Because production is constant returns to scale, total cost is proportional to output, so the multiplier $\lambda_{ij}$ is the nominal marginal cost:
\[
\lambda_{ij} = \text{mc}_{ij}.
\]

Solving for the exact unit-cost function yields
\begin{equation}
\text{mc}_{ij}
=
\frac{1}{\varphi_{ij}}
\left( \frac{w}{\beta_j} \right)^{\beta_j}
\left( \frac{P_j^M}{1-\beta_j} \right)^{1-\beta_j}.
\label{eq:app_exact_mc}
\end{equation}
Equivalently,
\[
\text{mc}_{ij}
=
\frac{w^{\beta_j} (P_j^M)^{1-\beta_j}}
{\varphi_{ij} \beta_j^{\beta_j} (1-\beta_j)^{1-\beta_j}}.
\]
Ignoring the sector-specific constant $\beta_j^{-\beta_j} (1-\beta_j)^{-(1-\beta_j)}$, the simpler proportionality used in the main text is
\begin{equation}
\text{mc}_{ij} \propto \varphi_{ij}^{-1} w^{\beta_j} (P_j^M)^{1-\beta_j}.
\label{eq:app_mc_level}
\end{equation}

Taking logs of \eqref{eq:app_exact_mc} or \eqref{eq:app_mc_level} gives
\begin{equation}
d \ln \text{mc}_{ij}
=
- d \ln \varphi_{ij}
+ \beta_j d \ln w
+ (1-\beta_j) d \ln P_j^M.
\label{eq:app_mc_growth}
\end{equation}

\subsection{Input-output propagation of Chinese input shocks}

Let $\mathbf{p}^M$ be the vector of sectoral log input prices, and let $A$ denote the domestic input-coefficient matrix. A first-order approximation of the network structure is
\[
\mathbf{p}^M = A \mathbf{p}^M + \Omega^{CHN} \mathbf{p}^{CHN},
\]
where $\mathbf{p}^{CHN}$ is the vector of Chinese input prices and $\Omega^{CHN}$ collects Chinese sourcing weights.

Rearranging,
\[
(I-A)\mathbf{p}^M = \Omega^{CHN} \mathbf{p}^{CHN},
\qquad
\mathbf{p}^M = (I-A)^{-1} \Omega^{CHN} \mathbf{p}^{CHN}.
\]
Let $L \equiv (I-A)^{-1}$ be the Leontief inverse. For sector $j$,
\begin{equation}
d \ln P_j^M
\approx
\sum_k l_{jk} \omega_k^{CHN} d \ln P_k^{CHN},
\label{eq:app_network_pass}
\end{equation}
where $l_{jk}$ is the $(j,k)$ entry of $L$ and $\omega_k^{CHN}$ is the relevant Chinese sourcing weight.

Substituting \eqref{eq:app_network_pass} into \eqref{eq:app_mc_growth} gives
\begin{equation}
d \ln \text{mc}_{ij}
\approx
- d \ln \varphi_{ij}
+ \beta_j d \ln w
+ (1-\beta_j) \sum_k l_{jk} \omega_k^{CHN} d \ln P_k^{CHN}.
\label{eq:app_mc_reduced}
\end{equation}

\subsection{Deriving the perceived elasticity step by step}

This is the core of the model. Since firms compete in quantities, the cleanest route is:
\begin{enumerate}
    \item write the inverse demand of firm $i$;
    \item differentiate it with respect to $q_{ij}$;
    \item track how a change in $q_{ij}$ moves $Q_j^D$, then $Q_j$, then $P_j$.
\end{enumerate}

\subsubsection{Step 1: inverse demand at the bottom nest}

From \eqref{eq:app_bottom_inverse_logs},
\[
\ln p_{ij}
=
\ln P_j^D
- \frac{1}{\rho_j} \ln q_{ij}
+ \frac{1}{\rho_j} \ln Q_j^D
+ \frac{1}{\rho_j} \ln \xi_i.
\]
Differentiate with respect to $\ln q_{ij}$:
\begin{equation}
\frac{d \ln p_{ij}}{d \ln q_{ij}}
=
\frac{d \ln P_j^D}{d \ln q_{ij}}
- \frac{1}{\rho_j}
+ \frac{1}{\rho_j} \frac{d \ln Q_j^D}{d \ln q_{ij}}.
\label{eq:app_step1}
\end{equation}
Using \eqref{eq:app_dlogQD_dlogq},
\begin{equation}
\frac{d \ln p_{ij}}{d \ln q_{ij}}
=
\frac{d \ln P_j^D}{d \ln q_{ij}}
- \frac{1-\tilde{s}_{ij}}{\rho_j}.
\label{eq:app_step1_simplified}
\end{equation}

\subsubsection{Step 2: how $P_j^D$ responds to $q_{ij}$}

From \eqref{eq:app_middle_inverse_diff},
\[
d \ln P_j^D
=
d \ln P_j
- \frac{1}{\nu_j} d \ln Q_j^D
+ \frac{1}{\nu_j} d \ln Q_j.
\]
Divide by $d \ln q_{ij}$:
\begin{equation}
\frac{d \ln P_j^D}{d \ln q_{ij}}
=
\frac{d \ln P_j}{d \ln q_{ij}}
- \frac{1}{\nu_j} \frac{d \ln Q_j^D}{d \ln q_{ij}}
+ \frac{1}{\nu_j} \frac{d \ln Q_j}{d \ln q_{ij}}.
\label{eq:app_step2}
\end{equation}
Using \eqref{eq:app_dlogQD_dlogq} and \eqref{eq:app_dlogQ_dlogq},
\begin{equation}
\frac{d \ln P_j^D}{d \ln q_{ij}}
=
\frac{d \ln P_j}{d \ln q_{ij}}
- \frac{\tilde{s}_{ij}}{\nu_j}
+ \frac{s_{ij}}{\nu_j}
=
\frac{d \ln P_j}{d \ln q_{ij}}
- \frac{\tilde{s}_{ij} - s_{ij}}{\nu_j}.
\label{eq:app_step2_simplified}
\end{equation}

\subsubsection{Step 3: how $P_j$ responds to $q_{ij}$}

From the sector-level inverse demand \eqref{eq:app_top_inverse},
\[
\ln P_j
=
\ln P
+ \frac{1}{\eta} \ln \alpha_j
- \frac{1}{\eta} \ln Q_j
+ \frac{1}{\eta} \ln C.
\]
Holding the aggregate objects $P$ and $C$ fixed from the viewpoint of an individual firm inside sector $j$,
\begin{equation}
\frac{d \ln P_j}{d \ln q_{ij}}
=
- \frac{1}{\eta} \frac{d \ln Q_j}{d \ln q_{ij}}
=
- \frac{s_{ij}}{\eta},
\label{eq:app_step3}
\end{equation}
where the last equality uses \eqref{eq:app_dlogQ_dlogq}.

\subsubsection{Step 4: combine the three pieces}

Insert \eqref{eq:app_step3} into \eqref{eq:app_step2_simplified}:
\[
\frac{d \ln P_j^D}{d \ln q_{ij}}
=
- \frac{\tilde{s}_{ij} - s_{ij}}{\nu_j}
- \frac{s_{ij}}{\eta}.
\]
Now plug this into \eqref{eq:app_step1_simplified}:
\[
\frac{d \ln p_{ij}}{d \ln q_{ij}}
=
- \frac{1-\tilde{s}_{ij}}{\rho_j}
- \frac{\tilde{s}_{ij} - s_{ij}}{\nu_j}
- \frac{s_{ij}}{\eta}.
\]
Therefore
\begin{equation}
- \frac{d \ln p_{ij}}{d \ln q_{ij}}
=
\frac{1-\tilde{s}_{ij}}{\rho_j}
+
\frac{\tilde{s}_{ij} - s_{ij}}{\nu_j}
+
\frac{s_{ij}}{\eta}.
\label{eq:app_inverse_demand_elasticity}
\end{equation}

Define the firm's \emph{perceived elasticity} under Cournot competition as
\begin{equation}
\frac{1}{\sigma_{ij}^C}
\equiv
- \frac{d \ln p_{ij}}{d \ln q_{ij}}.
\label{eq:app_sigma_definition}
\end{equation}
Equivalently, since price and quantity are locally inverse functions,
\[
\sigma_{ij}^C
=
- \frac{d \ln q_{ij}}{d \ln p_{ij}}.
\]
Combining \eqref{eq:app_inverse_demand_elasticity} and \eqref{eq:app_sigma_definition} yields
\begin{equation}
\sigma_{ij}^C
=
\left(
\frac{1-\tilde{s}_{ij}}{\rho_j}
+
\frac{\tilde{s}_{ij} - s_{ij}}{\nu_j}
+
\frac{s_{ij}}{\eta}
\right)^{-1}.
\label{eq:app_sigma_nested}
\end{equation}

\paragraph{Interpretation.}
Equation \eqref{eq:app_sigma_nested} is easiest to read term by term.
\begin{itemize}
    \item The term $\frac{1-\tilde{s}_{ij}}{\rho_j}$ comes from substitution across \emph{domestic varieties}. If the firm is small within the domestic segment, this force is strong.
    \item The term $\frac{\tilde{s}_{ij}-s_{ij}}{\nu_j}$ comes from substitution between the \emph{domestic composite} and the \emph{foreign composite} inside the sector.
    \item The term $\frac{s_{ij}}{\eta}$ comes from substitution across \emph{sectors} in final demand.
\end{itemize}
Because $\rho_j > \nu_j > \eta$, the inverse elasticities satisfy
\[
\frac{1}{\rho_j} < \frac{1}{\nu_j} < \frac{1}{\eta}.
\]
So as a firm becomes more important in broader and broader aggregates, the relevant residual demand it faces becomes less elastic.

\subsection{Cournot firm's FOC and the markup formula}

Firm $i$ chooses quantity $q_{ij}$ to maximize
\[
\pi_{ij}
=
\bigl(p_{ij}(q_{ij},q_{-ij}) - \text{mc}_{ij}\bigr) q_{ij} - f_{ij}.
\]
The first-order condition is
\begin{equation}
\frac{\partial \pi_{ij}}{\partial q_{ij}}
=
p_{ij}
+ q_{ij} \frac{\partial p_{ij}}{\partial q_{ij}}
- \text{mc}_{ij}
=
0.
\label{eq:app_cournot_foc}
\end{equation}

Divide \eqref{eq:app_cournot_foc} by $p_{ij}$:
\[
1
+ \frac{q_{ij}}{p_{ij}} \frac{\partial p_{ij}}{\partial q_{ij}}
- \frac{\text{mc}_{ij}}{p_{ij}}
=
0.
\]
Using the definition of the perceived elasticity,
\[
\frac{q_{ij}}{p_{ij}} \frac{\partial p_{ij}}{\partial q_{ij}}
=
\frac{d \ln p_{ij}}{d \ln q_{ij}}
=
- \frac{1}{\sigma_{ij}^C}.
\]
Hence
\[
1 - \frac{1}{\sigma_{ij}^C} = \frac{\text{mc}_{ij}}{p_{ij}}.
\]
Define the markup
\[
\mu_{ij} \equiv \frac{p_{ij}}{\text{mc}_{ij}}.
\]
Then
\begin{equation}
\frac{1}{\mu_{ij}}
=
1 - \frac{1}{\sigma_{ij}^C},
\qquad
\mu_{ij}
=
\frac{1}{1-\frac{1}{\sigma_{ij}^C}}
=
\frac{\sigma_{ij}^C}{\sigma_{ij}^C-1}.
\label{eq:app_markup_sigma}
\end{equation}

Substituting \eqref{eq:app_sigma_nested} into \eqref{eq:app_markup_sigma},
\begin{equation}
\frac{1}{\mu_{ij}}
=
1
-
\left[
\frac{1-\tilde{s}_{ij}}{\rho_j}
+
\frac{\tilde{s}_{ij} - s_{ij}}{\nu_j}
+
\frac{s_{ij}}{\eta}
\right].
\label{eq:app_inverse_markup_unsimplified}
\end{equation}
This is already a useful formula: a larger firm faces a less elastic residual demand, so $1/\mu_{ij}$ is lower and $\mu_{ij}$ is higher.

Now use $s_{ij} = \lambda_j^D \tilde{s}_{ij}$, so $\tilde{s}_{ij} = s_{ij}/\lambda_j^D$. Then \eqref{eq:app_inverse_markup_unsimplified} becomes
\begin{align}
\frac{1}{\mu_{ij}}
&=
1 - \frac{1}{\rho_j}
+ \frac{\tilde{s}_{ij}}{\rho_j}
- \frac{\tilde{s}_{ij}}{\nu_j}
+ \frac{s_{ij}}{\nu_j}
- \frac{s_{ij}}{\eta}
\notag \\
&=
\left( 1 - \frac{1}{\rho_j} \right)
-
\left[
\frac{1}{\lambda_j^D}
\left( \frac{1}{\nu_j} - \frac{1}{\rho_j} \right)
+
\left( \frac{1}{\eta} - \frac{1}{\nu_j} \right)
\right] s_{ij}.
\label{eq:app_inverse_markup_linear}
\end{align}

\paragraph{Important notation point.}
Equation \eqref{eq:app_inverse_markup_linear} is a linear equation for the \emph{inverse markup} $1/\mu_{ij}$, not for $\mu_{ij}$ itself. Therefore the corresponding markup is
\begin{equation}
\mu_{ij}
=
\frac{1}{
\left( 1 - \frac{1}{\rho_j} \right)
-
\left[
\frac{1}{\lambda_j^D}
\left( \frac{1}{\nu_j} - \frac{1}{\rho_j} \right)
+
\left( \frac{1}{\eta} - \frac{1}{\nu_j} \right)
\right] s_{ij}
}.
\label{eq:app_markup_full}
\end{equation}
In the main text, I suppress the distinction between $\lambda_j^D$ and $\Lambda_j^D$ for expositional simplicity; under that notation, \eqref{eq:app_inverse_markup_linear} becomes exactly the linear markup condition reported there, except that it should strictly be read as an inverse-markup condition.

\subsection{Sectoral markup and concentration}

Define the sectoral markup as the share-weighted harmonic mean of firm markups:
\begin{equation}
\mu_j
\equiv
\left(
\sum_{i=1}^{N_j} \frac{s_{ij}}{\mu_{ij}}
\right)^{-1}.
\label{eq:app_sector_harmonic}
\end{equation}
Taking inverses,
\[
\frac{1}{\mu_j}
=
\sum_{i=1}^{N_j} s_{ij} \frac{1}{\mu_{ij}}.
\]
Now substitute \eqref{eq:app_inverse_markup_linear}:
\begin{align}
\frac{1}{\mu_j}
&=
\sum_{i=1}^{N_j}
s_{ij}
\left[
\left( 1 - \frac{1}{\rho_j} \right)
-
\left(
\frac{1}{\lambda_j^D}
\left( \frac{1}{\nu_j} - \frac{1}{\rho_j} \right)
+
\left( \frac{1}{\eta} - \frac{1}{\nu_j} \right)
\right) s_{ij}
\right]
\notag \\
&=
\left( 1 - \frac{1}{\rho_j} \right) \sum_{i=1}^{N_j} s_{ij}
-
\left[
\frac{1}{\lambda_j^D}
\left( \frac{1}{\nu_j} - \frac{1}{\rho_j} \right)
+
\left( \frac{1}{\eta} - \frac{1}{\nu_j} \right)
\right]
\sum_{i=1}^{N_j} s_{ij}^2.
\label{eq:app_sector_inverse_step}
\end{align}
Since $\sum_i s_{ij} = 1$, and defining the Herfindahl index
\[
\text{HHI}_j \equiv \sum_{i=1}^{N_j} s_{ij}^2,
\]
we obtain
\begin{equation}
\frac{1}{\mu_j}
=
\left( 1 - \frac{1}{\rho_j} \right)
-
\left[
\frac{1}{\lambda_j^D}
\left( \frac{1}{\nu_j} - \frac{1}{\rho_j} \right)
+
\left( \frac{1}{\eta} - \frac{1}{\nu_j} \right)
\right]
\text{HHI}_j.
\label{eq:app_sector_inverse_markup}
\end{equation}
Therefore
\begin{equation}
\mu_j
=
\frac{1}{
\left( 1 - \frac{1}{\rho_j} \right)
-
\left[
\frac{1}{\lambda_j^D}
\left( \frac{1}{\nu_j} - \frac{1}{\rho_j} \right)
+
\left( \frac{1}{\eta} - \frac{1}{\nu_j} \right)
\right]
\text{HHI}_j
}.
\label{eq:app_sector_markup}
\end{equation}

Again, the \emph{linear} object is $1/\mu_j$, not $\mu_j$.

\subsubsection{Why concentration matters}

Let
\[
A_j \equiv 1 - \frac{1}{\rho_j},
\qquad
B_j \equiv
\frac{1}{\lambda_j^D}
\left( \frac{1}{\nu_j} - \frac{1}{\rho_j} \right)
+
\left( \frac{1}{\eta} - \frac{1}{\nu_j} \right).
\]
Then
\[
\frac{1}{\mu_j} = A_j - B_j \text{HHI}_j,
\qquad
\mu_j = \frac{1}{A_j - B_j \text{HHI}_j}.
\]
Differentiating,
\[
\frac{\partial \mu_j}{\partial \text{HHI}_j}
=
\frac{B_j}{(A_j - B_j \text{HHI}_j)^2}.
\]
Under $\rho_j > \nu_j > \eta$, we have
\[
\frac{1}{\nu_j} - \frac{1}{\rho_j} > 0,
\qquad
\frac{1}{\eta} - \frac{1}{\nu_j} > 0,
\]
so $B_j > 0$. Hence
\begin{equation}
\frac{\partial \mu_j}{\partial \text{HHI}_j} > 0.
\label{eq:app_markup_concentration_positive}
\end{equation}
Sectoral markups rise with concentration because more expenditure gets concentrated on firms that face less elastic residual demand.

\subsection{Dynamic extension and endogenous exit}

Let $x_{ijt}$ denote the state of firm $i$ in sector $j$ at date $t$. This state includes productivity, demand shifters, lagged market share, and the output- and input-trade states $(\tau_{jt}^O,\tau_{jt}^I)$.

If the firm stays active, its value function is
\begin{equation}
V(x_{ijt})
=
\max
\left\{
0,
\pi(x_{ijt})
+ \beta \mathbb{E}\left[ V(x_{ij,t+1}) \mid x_{ijt} \right]
\right\},
\label{eq:app_bellman}
\end{equation}
where $\beta \in (0,1)$ is the discount factor.

Define the continuation surplus
\[
W(x_{ijt})
\equiv
\pi(x_{ijt})
+ \beta \mathbb{E}\left[ V(x_{ij,t+1}) \mid x_{ijt} \right].
\]
The optimal exit rule is
\begin{equation}
\text{Exit}_{ijt} = 1
\quad \Longleftrightarrow \quad
W(x_{ijt}) < 0.
\label{eq:app_exit_rule}
\end{equation}
Equivalently, if current profitability can be summarized by a latent scalar $\theta_{ijt}$, there exists a cutoff function $\theta^*(x_{ijt})$ such that
\begin{equation}
\text{Exit}_{ijt} = 1
\quad \Longleftrightarrow \quad
\theta_{ijt} < \theta^*(x_{ijt}).
\label{eq:app_exit_threshold}
\end{equation}

\paragraph{Why this matters empirically.}
Trade shocks affect both current profits and future continuation values. A rise in final-good import competition may reduce current operating profits; a favorable input-supply shock may reduce marginal cost and increase the value of staying active. Therefore the sample of surviving firms is not random. Observed firm-level markup changes mix:
\begin{itemize}
    \item genuine within-incumbent adjustment;
    \item reallocation across surviving incumbents;
    \item selection induced by endogenous exit.
\end{itemize}
This is the reason the empirical section later introduces a selection-correction strategy.
.

\section{Theoretical Appendix}
\label{app:theory}

This appendix solves the model in Section 2 step by step. I write the algebra in a way that should be readable for a good undergraduate student who is already comfortable with constrained optimization, CES demand systems, and the idea of a markup as price over marginal cost.

The main goals are:
\begin{enumerate}
    \item derive the demand system at each CES nest from first principles;
    \item derive the inverse-demand system because firms compete in quantities;
    \item show carefully how the perceived elasticity is built from the three substitution margins in the model;
    \item derive the Cournot markup formula and the sectoral markup aggregator;
    \item clarify one notation point: the linear expression in market shares is the \emph{inverse markup} $1/\mu_{ij}$, not the markup $\mu_{ij}$ itself.
\end{enumerate}

\subsection{A useful CES template}

Before solving each layer separately, it is useful to record one generic CES result that we will reuse.

Suppose a composite $X$ is built from a collection of goods $\{x_n\}_{n \in \mathcal{N}}$:
\[
X =
\left(
\sum_{n \in \mathcal{N}}
a_n^{\frac{1}{\sigma}} x_n^{\frac{\sigma-1}{\sigma}}
\right)^{\frac{\sigma}{\sigma-1}},
\qquad \sigma > 1.
\]
Given prices $\{p_n\}$, the dual expenditure-minimization problem is
\[
\min_{\{x_n\}} \sum_{n \in \mathcal{N}} p_n x_n
\qquad \text{s.t.} \qquad
X \geq \bar{X}.
\]
The Lagrangian is
\[
\mathcal{L}
=
\sum_{n \in \mathcal{N}} p_n x_n
+
\lambda
\left[
\bar{X}
-
\left(
\sum_{n \in \mathcal{N}}
a_n^{\frac{1}{\sigma}} x_n^{\frac{\sigma-1}{\sigma}}
\right)^{\frac{\sigma}{\sigma-1}}
\right].
\]

Let
\[
M \equiv \sum_{n \in \mathcal{N}} a_n^{\frac{1}{\sigma}} x_n^{\frac{\sigma-1}{\sigma}},
\qquad
X = M^{\frac{\sigma}{\sigma-1}}.
\]
Then
\[
\frac{\partial X}{\partial x_n}
=
\frac{\sigma}{\sigma-1}
M^{\frac{1}{\sigma-1}}
\cdot
a_n^{\frac{1}{\sigma}}
\cdot
\frac{\sigma-1}{\sigma} x_n^{-\frac{1}{\sigma}}
=
X^{\frac{1}{\sigma}} a_n^{\frac{1}{\sigma}} x_n^{-\frac{1}{\sigma}}.
\]
Therefore the first-order condition for each $n$ is
\begin{equation}
p_n = \lambda X^{\frac{1}{\sigma}} a_n^{\frac{1}{\sigma}} x_n^{-\frac{1}{\sigma}}.
\label{eq:app_generic_ces_foc}
\end{equation}

Take two goods $n$ and $m$. Dividing the two FOCs gives
\[
\frac{p_n}{p_m}
=
\left( \frac{a_n}{a_m} \right)^{\frac{1}{\sigma}}
\left( \frac{x_n}{x_m} \right)^{-\frac{1}{\sigma}}
\Longrightarrow
\frac{x_n}{x_m}
=
\frac{a_n}{a_m}
\left( \frac{p_n}{p_m} \right)^{-\sigma}.
\]
This is the familiar CES relative-demand equation. Solving completely gives
\begin{equation}
x_n = a_n \left( \frac{p_n}{P_X} \right)^{-\sigma} X,
\label{eq:app_generic_hicksian}
\end{equation}
where the CES price index is
\begin{equation}
P_X =
\left(
\sum_{n \in \mathcal{N}}
a_n p_n^{1-\sigma}
\right)^{\frac{1}{1-\sigma}}.
\label{eq:app_generic_price_index}
\end{equation}

Finally, rearranging \eqref{eq:app_generic_hicksian} gives the inverse-demand form
\begin{equation}
p_n
=
P_X
\left( \frac{x_n}{a_n X} \right)^{-\frac{1}{\sigma}}
=
a_n^{\frac{1}{\sigma}} P_X \left( \frac{x_n}{X} \right)^{-\frac{1}{\sigma}}.
\label{eq:app_generic_inverse_demand}
\end{equation}

We now apply the same steps to each nest in the model.

\subsection{Top nest: allocation across sectors}

The representative household values sectoral composites according to
\[
C =
\left(
\int_0^1
\alpha_j^{\frac{1}{\eta}} Q_j^{\frac{\eta-1}{\eta}} dj
\right)^{\frac{\eta}{\eta-1}},
\qquad
\eta > 1,
\qquad
\int_0^1 \alpha_j dj = 1.
\]

\subsubsection{Expenditure minimization problem}

Fix a target final consumption level $\bar{C}$. The household solves
\[
\min_{\{Q_j\}_{j \in [0,1]}}
\int_0^1 P_j Q_j dj
\qquad \text{s.t.} \qquad
\left(
\int_0^1
\alpha_j^{\frac{1}{\eta}} Q_j^{\frac{\eta-1}{\eta}} dj
\right)^{\frac{\eta}{\eta-1}}
\geq \bar{C}.
\]

The Lagrangian is
\[
\mathcal{L}
=
\int_0^1 P_j Q_j dj
+
\lambda
\left[
\bar{C}
-
\left(
\int_0^1
\alpha_j^{\frac{1}{\eta}} Q_j^{\frac{\eta-1}{\eta}} dj
\right)^{\frac{\eta}{\eta-1}}
\right].
\]

Using the generic CES template with a continuum of goods, the FOC for sector $j$ is
\begin{equation}
P_j = \lambda C^{\frac{1}{\eta}} \alpha_j^{\frac{1}{\eta}} Q_j^{-\frac{1}{\eta}}.
\label{eq:app_top_foc}
\end{equation}
Take two sectors $j$ and $\ell$. Dividing their FOCs gives
\[
\frac{Q_j}{Q_\ell}
=
\frac{\alpha_j}{\alpha_\ell}
\left( \frac{P_j}{P_\ell} \right)^{-\eta}.
\]
Hence sectoral demand is
\begin{equation}
Q_j = \alpha_j \left( \frac{P_j}{P} \right)^{-\eta} C,
\label{eq:app_top_demand}
\end{equation}
where
\begin{equation}
P =
\left(
\int_0^1 \alpha_j P_j^{1-\eta} dj
\right)^{\frac{1}{1-\eta}}
\label{eq:app_top_price_index}
\end{equation}
is the final CES price index.

\subsubsection{Sectoral expenditure}

Multiply \eqref{eq:app_top_demand} by $P_j$:
\[
D_j \equiv P_j Q_j = \alpha_j \left( \frac{P_j}{P} \right)^{1-\eta} P C.
\]
Using $Y = PC$, we obtain
\begin{equation}
D_j = \alpha_j \left( \frac{P_j}{P} \right)^{1-\eta} Y.
\label{eq:app_sector_expenditure}
\end{equation}
This is the sector-level expenditure expression reported in the main text.

\subsubsection{Inverse demand at the sector level}

Rearranging \eqref{eq:app_top_demand} gives
\begin{equation}
P_j = P \, \alpha_j^{\frac{1}{\eta}} \left( \frac{Q_j}{C} \right)^{-\frac{1}{\eta}}.
\label{eq:app_top_inverse}
\end{equation}
Taking logs,
\[
\ln P_j = \ln P + \frac{1}{\eta} \ln \alpha_j - \frac{1}{\eta} \ln Q_j + \frac{1}{\eta} \ln C.
\]
So, holding $P$, $\alpha_j$, and $C$ fixed,
\begin{equation}
\frac{d \ln P_j}{d \ln Q_j} = - \frac{1}{\eta}.
\label{eq:app_top_inverse_elasticity}
\end{equation}
This is the inverse-demand elasticity at the top nest.

\subsection{Middle nest: domestic versus foreign supply inside sector}

Within sector $j$, the sectoral composite is
\[
Q_j =
\left[
(\Lambda_j^D)^{\frac{1}{\nu_j}} (Q_j^D)^{\frac{\nu_j-1}{\nu_j}}
+
(1-\Lambda_j^D)^{\frac{1}{\nu_j}} (Q_j^F)^{\frac{\nu_j-1}{\nu_j}}
\right]^{\frac{\nu_j}{\nu_j-1}},
\qquad \nu_j > 1.
\]

\subsubsection{Expenditure minimization problem}

For a given target $\bar{Q}_j$, minimize expenditure on the domestic and foreign composites:
\[
\min_{Q_j^D,Q_j^F}
P_j^D Q_j^D + P_j^F Q_j^F
\qquad \text{s.t.} \qquad
Q_j \geq \bar{Q}_j.
\]
The Lagrangian is
\[
\mathcal{L}
=
P_j^D Q_j^D + P_j^F Q_j^F
+
\lambda_j
\left[
\bar{Q}_j
-
\left[
(\Lambda_j^D)^{\frac{1}{\nu_j}} (Q_j^D)^{\frac{\nu_j-1}{\nu_j}}
+
(1-\Lambda_j^D)^{\frac{1}{\nu_j}} (Q_j^F)^{\frac{\nu_j-1}{\nu_j}}
\right]^{\frac{\nu_j}{\nu_j-1}}
\right].
\]

The FOCs are
\begin{equation}
P_j^D
=
\lambda_j Q_j^{\frac{1}{\nu_j}}
(\Lambda_j^D)^{\frac{1}{\nu_j}}
(Q_j^D)^{-\frac{1}{\nu_j}},
\label{eq:app_middle_foc_D}
\end{equation}
\begin{equation}
P_j^F
=
\lambda_j Q_j^{\frac{1}{\nu_j}}
(1-\Lambda_j^D)^{\frac{1}{\nu_j}}
(Q_j^F)^{-\frac{1}{\nu_j}}.
\label{eq:app_middle_foc_F}
\end{equation}

Dividing \eqref{eq:app_middle_foc_D} by \eqref{eq:app_middle_foc_F} gives
\[
\frac{Q_j^D}{Q_j^F}
=
\frac{\Lambda_j^D}{1-\Lambda_j^D}
\left( \frac{P_j^D}{P_j^F} \right)^{-\nu_j}.
\]
Therefore the Hicksian demands are
\begin{equation}
Q_j^D = \Lambda_j^D \left( \frac{P_j^D}{P_j} \right)^{-\nu_j} Q_j,
\label{eq:app_middle_demand_D}
\end{equation}
\begin{equation}
Q_j^F = (1-\Lambda_j^D) \left( \frac{P_j^F}{P_j} \right)^{-\nu_j} Q_j,
\label{eq:app_middle_demand_F}
\end{equation}
with price index
\begin{equation}
P_j =
\left[
\Lambda_j^D (P_j^D)^{1-\nu_j}
+
(1-\Lambda_j^D) (P_j^F)^{1-\nu_j}
\right]^{\frac{1}{1-\nu_j}}.
\label{eq:app_middle_price_index}
\end{equation}

\subsubsection{A useful notation point}

The CES weight $\Lambda_j^D$ is not, in general, the same thing as the realized domestic expenditure share. The realized share is
\begin{equation}
\lambda_j^D
\equiv
\frac{P_j^D Q_j^D}{P_j Q_j}
=
\Lambda_j^D \left( \frac{P_j^D}{P_j} \right)^{1-\nu_j}.
\label{eq:app_realized_domestic_share}
\end{equation}
In the main text, I often use $\Lambda_j^D$ as reduced-form shorthand for the domestic expenditure share. For the derivations below, it is cleaner to keep the distinction:
\begin{itemize}
    \item $\Lambda_j^D$: CES weight in preferences;
    \item $\lambda_j^D$: realized domestic expenditure share.
\end{itemize}

\subsubsection{Inverse demand for the domestic composite}

Rearranging \eqref{eq:app_middle_demand_D} gives
\begin{equation}
P_j^D
=
P_j
\left( \frac{Q_j^D}{\Lambda_j^D Q_j} \right)^{-\frac{1}{\nu_j}}.
\label{eq:app_middle_inverse_D}
\end{equation}
Taking logs,
\[
\ln P_j^D
=
\ln P_j
- \frac{1}{\nu_j} \ln Q_j^D
+ \frac{1}{\nu_j} \ln Q_j
+ \frac{1}{\nu_j} \ln \Lambda_j^D.
\]
Hence, when $\Lambda_j^D$ is treated as given,
\begin{equation}
d \ln P_j^D
=
d \ln P_j
- \frac{1}{\nu_j} d \ln Q_j^D
+ \frac{1}{\nu_j} d \ln Q_j.
\label{eq:app_middle_inverse_diff}
\end{equation}

\subsection{Bottom nest: domestic varieties inside the domestic composite}

Within the domestic nest of sector $j$, the domestic composite is
\[
Q_j^D =
\left(
\sum_{i=1}^{N_j}
\xi_i^{\frac{1}{\rho_j}} q_{ij}^{\frac{\rho_j-1}{\rho_j}}
\right)^{\frac{\rho_j}{\rho_j-1}},
\qquad \rho_j > 1.
\]

\subsubsection{Expenditure minimization problem}

For a given target $\bar{Q}_j^D$, minimize spending across domestic varieties:
\[
\min_{\{q_{ij}\}_{i=1}^{N_j}}
\sum_{i=1}^{N_j} p_{ij} q_{ij}
\qquad \text{s.t.} \qquad
Q_j^D \geq \bar{Q}_j^D.
\]
The Lagrangian is
\[
\mathcal{L}
=
\sum_{i=1}^{N_j} p_{ij} q_{ij}
+
\kappa_j
\left[
\bar{Q}_j^D
-
\left(
\sum_{i=1}^{N_j}
\xi_i^{\frac{1}{\rho_j}} q_{ij}^{\frac{\rho_j-1}{\rho_j}}
\right)^{\frac{\rho_j}{\rho_j-1}}
\right].
\]

The FOC for variety $i$ is
\begin{equation}
p_{ij}
=
\kappa_j (Q_j^D)^{\frac{1}{\rho_j}}
\xi_i^{\frac{1}{\rho_j}}
q_{ij}^{-\frac{1}{\rho_j}}.
\label{eq:app_bottom_foc}
\end{equation}
Take two varieties $i$ and $h$ and divide the FOCs:
\[
\frac{q_{ij}}{q_{hj}}
=
\frac{\xi_i}{\xi_h}
\left( \frac{p_{ij}}{p_{hj}} \right)^{-\rho_j}.
\]
Therefore the Hicksian demands are
\begin{equation}
q_{ij}
=
\xi_i \left( \frac{p_{ij}}{P_j^D} \right)^{-\rho_j} Q_j^D,
\label{eq:app_bottom_demand}
\end{equation}
where
\begin{equation}
P_j^D
=
\left(
\sum_{i=1}^{N_j}
\xi_i p_{ij}^{1-\rho_j}
\right)^{\frac{1}{1-\rho_j}}
\label{eq:app_bottom_price_index}
\end{equation}
is the CES price index for the domestic composite.

\subsubsection{Inverse demand for a domestic variety}

Rearranging \eqref{eq:app_bottom_demand},
\begin{equation}
p_{ij}
=
P_j^D
\left( \frac{q_{ij}}{\xi_i Q_j^D} \right)^{-\frac{1}{\rho_j}}
=
\xi_i^{\frac{1}{\rho_j}} P_j^D \left( \frac{q_{ij}}{Q_j^D} \right)^{-\frac{1}{\rho_j}}.
\label{eq:app_bottom_inverse}
\end{equation}
Taking logs gives
\begin{equation}
\ln p_{ij}
=
\ln P_j^D
- \frac{1}{\rho_j} \ln q_{ij}
+ \frac{1}{\rho_j} \ln Q_j^D
+ \frac{1}{\rho_j} \ln \xi_i.
\label{eq:app_bottom_inverse_logs}
\end{equation}

\subsubsection{Revenue shares}

Multiply \eqref{eq:app_bottom_demand} by $p_{ij}$ and divide by total domestic expenditure $P_j^D Q_j^D$:
\begin{equation}
\tilde{s}_{ij}
\equiv
\frac{p_{ij} q_{ij}}{P_j^D Q_j^D}
=
\xi_i \left( \frac{p_{ij}}{P_j^D} \right)^{1-\rho_j}.
\label{eq:app_within_share}
\end{equation}
This is firm $i$'s expenditure share within the \emph{domestic} segment of sector $j$.

Its share in \emph{total} sectoral expenditure is
\begin{equation}
s_{ij}
\equiv
\frac{p_{ij} q_{ij}}{P_j Q_j}
=
\frac{P_j^D Q_j^D}{P_j Q_j}
\cdot
\frac{p_{ij} q_{ij}}{P_j^D Q_j^D}
=
\lambda_j^D \tilde{s}_{ij}.
\label{eq:app_total_share}
\end{equation}
If one uses the reduced-form notation of the main text that absorbs relative-price terms into the domestic share, \eqref{eq:app_total_share} becomes $s_{ij} = \Lambda_j^D \tilde{s}_{ij}$.

\subsection{Useful derivative identities}

The next step is to compute how an increase in firm $i$'s own output $q_{ij}$ affects the relevant aggregates $Q_j^D$ and $Q_j$. These derivatives are what allow the three CES nests to show up inside the perceived elasticity.

\subsubsection{How a change in $q_{ij}$ affects the domestic composite $Q_j^D$}

Differentiate the domestic aggregator with respect to $q_{ij}$:
\[
\frac{\partial Q_j^D}{\partial q_{ij}}
=
(Q_j^D)^{\frac{1}{\rho_j}}
\xi_i^{\frac{1}{\rho_j}}
q_{ij}^{-\frac{1}{\rho_j}}.
\]
Multiply by $q_{ij}/Q_j^D$:
\[
\frac{\partial \ln Q_j^D}{\partial \ln q_{ij}}
=
\frac{q_{ij}}{Q_j^D}
\frac{\partial Q_j^D}{\partial q_{ij}}
=
\xi_i^{\frac{1}{\rho_j}} q_{ij}^{\frac{\rho_j-1}{\rho_j}} (Q_j^D)^{-\frac{\rho_j-1}{\rho_j}}.
\]
Now use the inverse-demand equation \eqref{eq:app_bottom_inverse}:
\[
\frac{p_{ij} q_{ij}}{P_j^D Q_j^D}
=
\xi_i^{\frac{1}{\rho_j}} q_{ij}^{\frac{\rho_j-1}{\rho_j}} (Q_j^D)^{-\frac{\rho_j-1}{\rho_j}}
=
\tilde{s}_{ij}.
\]
Therefore
\begin{equation}
\frac{\partial \ln Q_j^D}{\partial \ln q_{ij}} = \tilde{s}_{ij}.
\label{eq:app_dlogQD_dlogq}
\end{equation}
Intuition: if firm $i$ is large within the domestic segment, increasing its own output has a larger effect on the domestic composite.

\subsubsection{How a change in $Q_j^D$ affects the sectoral composite $Q_j$}

Differentiate the middle CES aggregator with respect to $Q_j^D$:
\[
\frac{\partial Q_j}{\partial Q_j^D}
=
Q_j^{\frac{1}{\nu_j}}
(\Lambda_j^D)^{\frac{1}{\nu_j}}
(Q_j^D)^{-\frac{1}{\nu_j}}.
\]
Multiply by $Q_j^D/Q_j$:
\[
\frac{\partial \ln Q_j}{\partial \ln Q_j^D}
=
(\Lambda_j^D)^{\frac{1}{\nu_j}} (Q_j^D)^{\frac{\nu_j-1}{\nu_j}} Q_j^{-\frac{\nu_j-1}{\nu_j}}.
\]
Using \eqref{eq:app_middle_inverse_D}, this expression is exactly the realized domestic expenditure share:
\begin{equation}
\frac{\partial \ln Q_j}{\partial \ln Q_j^D}
=
\frac{P_j^D Q_j^D}{P_j Q_j}
=
\lambda_j^D.
\label{eq:app_dlogQ_dlogQD}
\end{equation}

Combine \eqref{eq:app_dlogQD_dlogq} and \eqref{eq:app_dlogQ_dlogQD}:
\begin{equation}
\frac{\partial \ln Q_j}{\partial \ln q_{ij}}
=
\frac{\partial \ln Q_j}{\partial \ln Q_j^D}
\cdot
\frac{\partial \ln Q_j^D}{\partial \ln q_{ij}}
=
\lambda_j^D \tilde{s}_{ij}
=
s_{ij}.
\label{eq:app_dlogQ_dlogq}
\end{equation}
So the effect of firm $i$'s quantity on the whole sectoral composite is exactly its total sectoral expenditure share.

\subsection{Cost minimization and marginal cost}

Domestic firm $i$ in sector $j$ produces with
\[
q_{ij} = \varphi_{ij} L_{ij}^{\beta_j} M_{ij}^{1-\beta_j},
\qquad 0 < \beta_j < 1.
\]
Given factor prices $w$ and $P_j^M$, the cost-minimization problem is
\[
\min_{L_{ij},M_{ij}}
w L_{ij} + P_j^M M_{ij}
\qquad \text{s.t.} \qquad
q_{ij} = \varphi_{ij} L_{ij}^{\beta_j} M_{ij}^{1-\beta_j}.
\]
The Lagrangian is
\[
\mathcal{L}
=
w L_{ij} + P_j^M M_{ij}
+
\lambda_{ij}
\left[
q_{ij} - \varphi_{ij} L_{ij}^{\beta_j} M_{ij}^{1-\beta_j}
\right].
\]

The FOCs are
\begin{equation}
w
=
\lambda_{ij} \varphi_{ij} \beta_j L_{ij}^{\beta_j-1} M_{ij}^{1-\beta_j},
\label{eq:app_cost_foc_L}
\end{equation}
\begin{equation}
P_j^M
=
\lambda_{ij} \varphi_{ij} (1-\beta_j) L_{ij}^{\beta_j} M_{ij}^{-\beta_j}.
\label{eq:app_cost_foc_M}
\end{equation}

Divide \eqref{eq:app_cost_foc_L} by \eqref{eq:app_cost_foc_M}:
\[
\frac{w}{P_j^M}
=
\frac{\beta_j}{1-\beta_j}
\frac{M_{ij}}{L_{ij}}.
\]
Hence the cost-minimizing input ratio is
\begin{equation}
\frac{M_{ij}}{L_{ij}}
=
\frac{1-\beta_j}{\beta_j} \frac{w}{P_j^M}.
\label{eq:app_input_ratio}
\end{equation}

Now multiply \eqref{eq:app_cost_foc_L} by $L_{ij}$ and \eqref{eq:app_cost_foc_M} by $M_{ij}$:
\[
w L_{ij}
=
\lambda_{ij} \beta_j q_{ij},
\qquad
P_j^M M_{ij}
=
\lambda_{ij} (1-\beta_j) q_{ij}.
\]
Adding the two equations gives
\[
w L_{ij} + P_j^M M_{ij} = \lambda_{ij} q_{ij}.
\]
Because production is constant returns to scale, total cost is proportional to output, so the multiplier $\lambda_{ij}$ is the nominal marginal cost:
\[
\lambda_{ij} = \text{mc}_{ij}.
\]

Solving for the exact unit-cost function yields
\begin{equation}
\text{mc}_{ij}
=
\frac{1}{\varphi_{ij}}
\left( \frac{w}{\beta_j} \right)^{\beta_j}
\left( \frac{P_j^M}{1-\beta_j} \right)^{1-\beta_j}.
\label{eq:app_exact_mc}
\end{equation}
Equivalently,
\[
\text{mc}_{ij}
=
\frac{w^{\beta_j} (P_j^M)^{1-\beta_j}}
{\varphi_{ij} \beta_j^{\beta_j} (1-\beta_j)^{1-\beta_j}}.
\]
Ignoring the sector-specific constant $\beta_j^{-\beta_j} (1-\beta_j)^{-(1-\beta_j)}$, the simpler proportionality used in the main text is
\begin{equation}
\text{mc}_{ij} \propto \varphi_{ij}^{-1} w^{\beta_j} (P_j^M)^{1-\beta_j}.
\label{eq:app_mc_level}
\end{equation}

Taking logs of \eqref{eq:app_exact_mc} or \eqref{eq:app_mc_level} gives
\begin{equation}
d \ln \text{mc}_{ij}
=
- d \ln \varphi_{ij}
+ \beta_j d \ln w
+ (1-\beta_j) d \ln P_j^M.
\label{eq:app_mc_growth}
\end{equation}

\subsection{Input-output propagation of Chinese input shocks}

Let $\mathbf{p}^M$ be the vector of sectoral log input prices, and let $A$ denote the domestic input-coefficient matrix. A first-order approximation of the network structure is
\[
\mathbf{p}^M = A \mathbf{p}^M + \Omega^{CHN} \mathbf{p}^{CHN},
\]
where $\mathbf{p}^{CHN}$ is the vector of Chinese input prices and $\Omega^{CHN}$ collects Chinese sourcing weights.

Rearranging,
\[
(I-A)\mathbf{p}^M = \Omega^{CHN} \mathbf{p}^{CHN},
\qquad
\mathbf{p}^M = (I-A)^{-1} \Omega^{CHN} \mathbf{p}^{CHN}.
\]
Let $L \equiv (I-A)^{-1}$ be the Leontief inverse. For sector $j$,
\begin{equation}
d \ln P_j^M
\approx
\sum_k l_{jk} \omega_k^{CHN} d \ln P_k^{CHN},
\label{eq:app_network_pass}
\end{equation}
where $l_{jk}$ is the $(j,k)$ entry of $L$ and $\omega_k^{CHN}$ is the relevant Chinese sourcing weight.

Substituting \eqref{eq:app_network_pass} into \eqref{eq:app_mc_growth} gives
\begin{equation}
d \ln \text{mc}_{ij}
\approx
- d \ln \varphi_{ij}
+ \beta_j d \ln w
+ (1-\beta_j) \sum_k l_{jk} \omega_k^{CHN} d \ln P_k^{CHN}.
\label{eq:app_mc_reduced}
\end{equation}

\subsection{Deriving the perceived elasticity step by step}

This is the core of the model. Since firms compete in quantities, the cleanest route is:
\begin{enumerate}
    \item write the inverse demand of firm $i$;
    \item differentiate it with respect to $q_{ij}$;
    \item track how a change in $q_{ij}$ moves $Q_j^D$, then $Q_j$, then $P_j$.
\end{enumerate}

\subsubsection{Step 1: inverse demand at the bottom nest}

From \eqref{eq:app_bottom_inverse_logs},
\[
\ln p_{ij}
=
\ln P_j^D
- \frac{1}{\rho_j} \ln q_{ij}
+ \frac{1}{\rho_j} \ln Q_j^D
+ \frac{1}{\rho_j} \ln \xi_i.
\]
Differentiate with respect to $\ln q_{ij}$:
\begin{equation}
\frac{d \ln p_{ij}}{d \ln q_{ij}}
=
\frac{d \ln P_j^D}{d \ln q_{ij}}
- \frac{1}{\rho_j}
+ \frac{1}{\rho_j} \frac{d \ln Q_j^D}{d \ln q_{ij}}.
\label{eq:app_step1}
\end{equation}
Using \eqref{eq:app_dlogQD_dlogq},
\begin{equation}
\frac{d \ln p_{ij}}{d \ln q_{ij}}
=
\frac{d \ln P_j^D}{d \ln q_{ij}}
- \frac{1-\tilde{s}_{ij}}{\rho_j}.
\label{eq:app_step1_simplified}
\end{equation}

\subsubsection{Step 2: how $P_j^D$ responds to $q_{ij}$}

From \eqref{eq:app_middle_inverse_diff},
\[
d \ln P_j^D
=
d \ln P_j
- \frac{1}{\nu_j} d \ln Q_j^D
+ \frac{1}{\nu_j} d \ln Q_j.
\]
Divide by $d \ln q_{ij}$:
\begin{equation}
\frac{d \ln P_j^D}{d \ln q_{ij}}
=
\frac{d \ln P_j}{d \ln q_{ij}}
- \frac{1}{\nu_j} \frac{d \ln Q_j^D}{d \ln q_{ij}}
+ \frac{1}{\nu_j} \frac{d \ln Q_j}{d \ln q_{ij}}.
\label{eq:app_step2}
\end{equation}
Using \eqref{eq:app_dlogQD_dlogq} and \eqref{eq:app_dlogQ_dlogq},
\begin{equation}
\frac{d \ln P_j^D}{d \ln q_{ij}}
=
\frac{d \ln P_j}{d \ln q_{ij}}
- \frac{\tilde{s}_{ij}}{\nu_j}
+ \frac{s_{ij}}{\nu_j}
=
\frac{d \ln P_j}{d \ln q_{ij}}
- \frac{\tilde{s}_{ij} - s_{ij}}{\nu_j}.
\label{eq:app_step2_simplified}
\end{equation}

\subsubsection{Step 3: how $P_j$ responds to $q_{ij}$}

From the sector-level inverse demand \eqref{eq:app_top_inverse},
\[
\ln P_j
=
\ln P
+ \frac{1}{\eta} \ln \alpha_j
- \frac{1}{\eta} \ln Q_j
+ \frac{1}{\eta} \ln C.
\]
Holding the aggregate objects $P$ and $C$ fixed from the viewpoint of an individual firm inside sector $j$,
\begin{equation}
\frac{d \ln P_j}{d \ln q_{ij}}
=
- \frac{1}{\eta} \frac{d \ln Q_j}{d \ln q_{ij}}
=
- \frac{s_{ij}}{\eta},
\label{eq:app_step3}
\end{equation}
where the last equality uses \eqref{eq:app_dlogQ_dlogq}.

\subsubsection{Step 4: combine the three pieces}

Insert \eqref{eq:app_step3} into \eqref{eq:app_step2_simplified}:
\[
\frac{d \ln P_j^D}{d \ln q_{ij}}
=
- \frac{\tilde{s}_{ij} - s_{ij}}{\nu_j}
- \frac{s_{ij}}{\eta}.
\]
Now plug this into \eqref{eq:app_step1_simplified}:
\[
\frac{d \ln p_{ij}}{d \ln q_{ij}}
=
- \frac{1-\tilde{s}_{ij}}{\rho_j}
- \frac{\tilde{s}_{ij} - s_{ij}}{\nu_j}
- \frac{s_{ij}}{\eta}.
\]
Therefore
\begin{equation}
- \frac{d \ln p_{ij}}{d \ln q_{ij}}
=
\frac{1-\tilde{s}_{ij}}{\rho_j}
+
\frac{\tilde{s}_{ij} - s_{ij}}{\nu_j}
+
\frac{s_{ij}}{\eta}.
\label{eq:app_inverse_demand_elasticity}
\end{equation}

Define the firm's \emph{perceived elasticity} under Cournot competition as
\begin{equation}
\frac{1}{\sigma_{ij}^C}
\equiv
- \frac{d \ln p_{ij}}{d \ln q_{ij}}.
\label{eq:app_sigma_definition}
\end{equation}
Equivalently, since price and quantity are locally inverse functions,
\[
\sigma_{ij}^C
=
- \frac{d \ln q_{ij}}{d \ln p_{ij}}.
\]
Combining \eqref{eq:app_inverse_demand_elasticity} and \eqref{eq:app_sigma_definition} yields
\begin{equation}
\sigma_{ij}^C
=
\left(
\frac{1-\tilde{s}_{ij}}{\rho_j}
+
\frac{\tilde{s}_{ij} - s_{ij}}{\nu_j}
+
\frac{s_{ij}}{\eta}
\right)^{-1}.
\label{eq:app_sigma_nested}
\end{equation}

\paragraph{Interpretation.}
Equation \eqref{eq:app_sigma_nested} is easiest to read term by term.
\begin{itemize}
    \item The term $\frac{1-\tilde{s}_{ij}}{\rho_j}$ comes from substitution across \emph{domestic varieties}. If the firm is small within the domestic segment, this force is strong.
    \item The term $\frac{\tilde{s}_{ij}-s_{ij}}{\nu_j}$ comes from substitution between the \emph{domestic composite} and the \emph{foreign composite} inside the sector.
    \item The term $\frac{s_{ij}}{\eta}$ comes from substitution across \emph{sectors} in final demand.
\end{itemize}
Because $\rho_j > \nu_j > \eta$, the inverse elasticities satisfy
\[
\frac{1}{\rho_j} < \frac{1}{\nu_j} < \frac{1}{\eta}.
\]
So as a firm becomes more important in broader and broader aggregates, the relevant residual demand it faces becomes less elastic.

\subsection{Cournot firm's FOC and the markup formula}

Firm $i$ chooses quantity $q_{ij}$ to maximize
\[
\pi_{ij}
=
\bigl(p_{ij}(q_{ij},q_{-ij}) - \text{mc}_{ij}\bigr) q_{ij} - f_{ij}.
\]
The first-order condition is
\begin{equation}
\frac{\partial \pi_{ij}}{\partial q_{ij}}
=
p_{ij}
+ q_{ij} \frac{\partial p_{ij}}{\partial q_{ij}}
- \text{mc}_{ij}
=
0.
\label{eq:app_cournot_foc}
\end{equation}

Divide \eqref{eq:app_cournot_foc} by $p_{ij}$:
\[
1
+ \frac{q_{ij}}{p_{ij}} \frac{\partial p_{ij}}{\partial q_{ij}}
- \frac{\text{mc}_{ij}}{p_{ij}}
=
0.
\]
Using the definition of the perceived elasticity,
\[
\frac{q_{ij}}{p_{ij}} \frac{\partial p_{ij}}{\partial q_{ij}}
=
\frac{d \ln p_{ij}}{d \ln q_{ij}}
=
- \frac{1}{\sigma_{ij}^C}.
\]
Hence
\[
1 - \frac{1}{\sigma_{ij}^C} = \frac{\text{mc}_{ij}}{p_{ij}}.
\]
Define the markup
\[
\mu_{ij} \equiv \frac{p_{ij}}{\text{mc}_{ij}}.
\]
Then
\begin{equation}
\frac{1}{\mu_{ij}}
=
1 - \frac{1}{\sigma_{ij}^C},
\qquad
\mu_{ij}
=
\frac{1}{1-\frac{1}{\sigma_{ij}^C}}
=
\frac{\sigma_{ij}^C}{\sigma_{ij}^C-1}.
\label{eq:app_markup_sigma}
\end{equation}

Substituting \eqref{eq:app_sigma_nested} into \eqref{eq:app_markup_sigma},
\begin{equation}
\frac{1}{\mu_{ij}}
=
1
-
\left[
\frac{1-\tilde{s}_{ij}}{\rho_j}
+
\frac{\tilde{s}_{ij} - s_{ij}}{\nu_j}
+
\frac{s_{ij}}{\eta}
\right].
\label{eq:app_inverse_markup_unsimplified}
\end{equation}
This is already a useful formula: a larger firm faces a less elastic residual demand, so $1/\mu_{ij}$ is lower and $\mu_{ij}$ is higher.

Now use $s_{ij} = \lambda_j^D \tilde{s}_{ij}$, so $\tilde{s}_{ij} = s_{ij}/\lambda_j^D$. Then \eqref{eq:app_inverse_markup_unsimplified} becomes
\begin{align}
\frac{1}{\mu_{ij}}
&=
1 - \frac{1}{\rho_j}
+ \frac{\tilde{s}_{ij}}{\rho_j}
- \frac{\tilde{s}_{ij}}{\nu_j}
+ \frac{s_{ij}}{\nu_j}
- \frac{s_{ij}}{\eta}
\notag \\
&=
\left( 1 - \frac{1}{\rho_j} \right)
-
\left[
\frac{1}{\lambda_j^D}
\left( \frac{1}{\nu_j} - \frac{1}{\rho_j} \right)
+
\left( \frac{1}{\eta} - \frac{1}{\nu_j} \right)
\right] s_{ij}.
\label{eq:app_inverse_markup_linear}
\end{align}

\paragraph{Important notation point.}
Equation \eqref{eq:app_inverse_markup_linear} is a linear equation for the \emph{inverse markup} $1/\mu_{ij}$, not for $\mu_{ij}$ itself. Therefore the corresponding markup is
\begin{equation}
\mu_{ij}
=
\frac{1}{
\left( 1 - \frac{1}{\rho_j} \right)
-
\left[
\frac{1}{\lambda_j^D}
\left( \frac{1}{\nu_j} - \frac{1}{\rho_j} \right)
+
\left( \frac{1}{\eta} - \frac{1}{\nu_j} \right)
\right] s_{ij}
}.
\label{eq:app_markup_full}
\end{equation}
In the main text, I suppress the distinction between $\lambda_j^D$ and $\Lambda_j^D$ for expositional simplicity; under that notation, \eqref{eq:app_inverse_markup_linear} becomes exactly the linear markup condition reported there, except that it should strictly be read as an inverse-markup condition.

\subsection{Sectoral markup and concentration}

Define the sectoral markup as the share-weighted harmonic mean of firm markups:
\begin{equation}
\mu_j
\equiv
\left(
\sum_{i=1}^{N_j} \frac{s_{ij}}{\mu_{ij}}
\right)^{-1}.
\label{eq:app_sector_harmonic}
\end{equation}
Taking inverses,
\[
\frac{1}{\mu_j}
=
\sum_{i=1}^{N_j} s_{ij} \frac{1}{\mu_{ij}}.
\]
Now substitute \eqref{eq:app_inverse_markup_linear}:
\begin{align}
\frac{1}{\mu_j}
&=
\sum_{i=1}^{N_j}
s_{ij}
\left[
\left( 1 - \frac{1}{\rho_j} \right)
-
\left(
\frac{1}{\lambda_j^D}
\left( \frac{1}{\nu_j} - \frac{1}{\rho_j} \right)
+
\left( \frac{1}{\eta} - \frac{1}{\nu_j} \right)
\right) s_{ij}
\right]
\notag \\
&=
\left( 1 - \frac{1}{\rho_j} \right) \sum_{i=1}^{N_j} s_{ij}
-
\left[
\frac{1}{\lambda_j^D}
\left( \frac{1}{\nu_j} - \frac{1}{\rho_j} \right)
+
\left( \frac{1}{\eta} - \frac{1}{\nu_j} \right)
\right]
\sum_{i=1}^{N_j} s_{ij}^2.
\label{eq:app_sector_inverse_step}
\end{align}
Since $\sum_i s_{ij} = 1$, and defining the Herfindahl index
\[
\text{HHI}_j \equiv \sum_{i=1}^{N_j} s_{ij}^2,
\]
we obtain
\begin{equation}
\frac{1}{\mu_j}
=
\left( 1 - \frac{1}{\rho_j} \right)
-
\left[
\frac{1}{\lambda_j^D}
\left( \frac{1}{\nu_j} - \frac{1}{\rho_j} \right)
+
\left( \frac{1}{\eta} - \frac{1}{\nu_j} \right)
\right]
\text{HHI}_j.
\label{eq:app_sector_inverse_markup}
\end{equation}
Therefore
\begin{equation}
\mu_j
=
\frac{1}{
\left( 1 - \frac{1}{\rho_j} \right)
-
\left[
\frac{1}{\lambda_j^D}
\left( \frac{1}{\nu_j} - \frac{1}{\rho_j} \right)
+
\left( \frac{1}{\eta} - \frac{1}{\nu_j} \right)
\right]
\text{HHI}_j
}.
\label{eq:app_sector_markup}
\end{equation}

Again, the \emph{linear} object is $1/\mu_j$, not $\mu_j$.

\subsubsection{Why concentration matters}

Let
\[
A_j \equiv 1 - \frac{1}{\rho_j},
\qquad
B_j \equiv
\frac{1}{\lambda_j^D}
\left( \frac{1}{\nu_j} - \frac{1}{\rho_j} \right)
+
\left( \frac{1}{\eta} - \frac{1}{\nu_j} \right).
\]
Then
\[
\frac{1}{\mu_j} = A_j - B_j \text{HHI}_j,
\qquad
\mu_j = \frac{1}{A_j - B_j \text{HHI}_j}.
\]
Differentiating,
\[
\frac{\partial \mu_j}{\partial \text{HHI}_j}
=
\frac{B_j}{(A_j - B_j \text{HHI}_j)^2}.
\]
Under $\rho_j > \nu_j > \eta$, we have
\[
\frac{1}{\nu_j} - \frac{1}{\rho_j} > 0,
\qquad
\frac{1}{\eta} - \frac{1}{\nu_j} > 0,
\]
so $B_j > 0$. Hence
\begin{equation}
\frac{\partial \mu_j}{\partial \text{HHI}_j} > 0.
\label{eq:app_markup_concentration_positive}
\end{equation}
Sectoral markups rise with concentration because more expenditure gets concentrated on firms that face less elastic residual demand.

\subsection{Dynamic extension and endogenous exit}

Let $x_{ijt}$ denote the state of firm $i$ in sector $j$ at date $t$. This state includes productivity, demand shifters, lagged market share, and the output- and input-trade states $(\tau_{jt}^O,\tau_{jt}^I)$.

If the firm stays active, its value function is
\begin{equation}
V(x_{ijt})
=
\max
\left\{
0,
\pi(x_{ijt})
+ \beta \mathbb{E}\left[ V(x_{ij,t+1}) \mid x_{ijt} \right]
\right\},
\label{eq:app_bellman}
\end{equation}
where $\beta \in (0,1)$ is the discount factor.

Define the continuation surplus
\[
W(x_{ijt})
\equiv
\pi(x_{ijt})
+ \beta \mathbb{E}\left[ V(x_{ij,t+1}) \mid x_{ijt} \right].
\]
The optimal exit rule is
\begin{equation}
\text{Exit}_{ijt} = 1
\quad \Longleftrightarrow \quad
W(x_{ijt}) < 0.
\label{eq:app_exit_rule}
\end{equation}
Equivalently, if current profitability can be summarized by a latent scalar $\theta_{ijt}$, there exists a cutoff function $\theta^*(x_{ijt})$ such that
\begin{equation}
\text{Exit}_{ijt} = 1
\quad \Longleftrightarrow \quad
\theta_{ijt} < \theta^*(x_{ijt}).
\label{eq:app_exit_threshold}
\end{equation}

\paragraph{Why this matters empirically.}
Trade shocks affect both current profits and future continuation values. A rise in final-good import competition may reduce current operating profits; a favorable input-supply shock may reduce marginal cost and increase the value of staying active. Therefore the sample of surviving firms is not random. Observed firm-level markup changes mix:
\begin{itemize}
    \item genuine within-incumbent adjustment;
    \item reallocation across surviving incumbents;
    \item selection induced by endogenous exit.
\end{itemize}
This is the reason the empirical section later introduces a selection-correction strategy.
